% Standard ACM conference proceedings: two-column sigconf layout.
% Compile with latexmk -pdf main.tex (BibTeX runs automatically).
\documentclass[sigconf]{acmart}

% Keep the class defaults for margins, typography, and bibliography.
\AtBeginDocument{%
  \providecommand\BibTeX{{Bib\TeX}}}

% DRAFT metadata. Replace with the exact ACM rights-form information
% before preparing the final publication version.
\setcopyright{none}
\acmYear{2026}
\copyrightyear{2026}
\acmConference[CONF '26]{Conference Name (Replace Before Submission)}%
  {Month DD--DD, 2026}{City, Country}
\acmBooktitle{Conference Name (CONF '26), Month DD--DD, 2026, City, Country}
\acmDOI{}
\acmISBN{}

% Shared preamble for the notation-aligned Sections III and IV.
% Load once before \begin{document}; theorem names are package-specific.
% acmart already loads amsmath, amssymb, amsthm and booktabs.
% Reuse its package options, theorem style and default typography.
\makeatletter
\@ifundefined{btdtheorem}{\newtheorem{btdtheorem}{Theorem}[section]}{}
\@ifundefined{btdlemma}{\newtheorem{btdlemma}[btdtheorem]{Lemma}}{}
\@ifundefined{btdproposition}{\newtheorem{btdproposition}[btdtheorem]{Proposition}}{}
\@ifundefined{btdcorollary}{\newtheorem{btdcorollary}[btdtheorem]{Corollary}}{}
\makeatother


\begin{document}

\title[Fuzzy PI Control for bittide]{Fuzzy PI Control for bittide: Nonlinear Analysis and Transient Performance}

\author{First Author}
\email{first.author@example.com}
\affiliation{%
  \institution{University or Organization}
  \city{City}
  \country{Country}
}

\author{Second Author}
\email{second.author@example.com}
\affiliation{%
  \institution{University or Organization}
  \city{City}
  \country{Country}
}

\renewcommand{\shortauthors}{First Author and Second Author}

\begin{abstract}
Finite receiving buffers couple synchronization performance to transient
resource use in bittide networks. We propose a decentralized incremental
Fuzzy PI controller that shapes oscillator corrections using occupancy
error and its update-to-update change. We specify its full inference law,
derive an implementation-consistent PI calibration and characterize its
central nonlinear map. For a synchronous zero-delay fluid reference,
an invariant-domain and Lyapunov analysis, combined with interval
enclosures and a contracting tail, certifies an all-time buffer-excursion
benefit for a prescribed initialization: below 17.315 frames for Fuzzy
versus approximately 23.404 for matched-gain PI. Continuous-phase Callisto
experiments separately evaluate delayed local execution. After equal-budget
selection from 24 candidates per family, 186 untouched tests show
mean paired disagreement-RMS improvements of 16.78\% and 7.85\% in two
bounded-random latency regimes. Matched-gain startup tests also show
smaller buffer excursions. These gains depend on the operating regime
and carry a common-frequency-shift cost; slower-random confirmation
delimits their applicability. The results connect an explicit nonlinear
mechanism to conditional headroom and synchronization benefits.
\end{abstract}


% Replace these example CCS terms using https://dl.acm.org/ccs .
\begin{CCSXML}
<ccs2012>
 <concept>
  <concept_id>10010147</concept_id>
  <concept_desc>Computing methodologies</concept_desc>
  <concept_significance>500</concept_significance>
 </concept>
</ccs2012>
\end{CCSXML}
\ccsdesc[500]{Computing methodologies}
\keywords{Distributed synchronization, bittide, fuzzy control, logical synchrony, cyber-physical systems}

\maketitle

% Introduction v5 is the active candidate. The author's earlier unfinished
% opening is preserved separately in author_notes/ and in the full backup.
% Adapted from introduction_EN_v5.md; research claims retained.
\section{Introduction}
\label{sec:introduction}

Distributed cyber-physical computation requires nodes running on independent clocks to coordinate their computations and communications. For applications that depend on reproducible execution, the relationships between communicated events must remain well defined as physical timing varies. This creates a fundamental implementation challenge: a coherent logical execution must be realized using oscillators, communication links and storage resources whose physical behaviour is neither identical nor unlimited.

Logical synchrony provides a framework for addressing this challenge without distributing a global physical clock. In bittide, continuous frame exchange and elastic receiving buffers connect the physical operation of neighbouring nodes to their logical execution. Each node adjusts its oscillator frequency using locally observed buffer occupancies. These buffers consequently serve two roles: they provide the feedback signal for synchronization and absorb the imbalance between incoming frames and local consumption. The storage available to accommodate that imbalance is finite~\cite{lall2024logicalSynchrony}.

The resulting control problem concerns the entire trajectory towards an operating point. Initial oscillator differences and changes in frame-delivery time can consume buffer headroom before the network settles or while it responds to a disturbance. A small eventual frequency disagreement does not by itself establish that the intervening buffer trajectory remains admissible. The objective is therefore to use local occupancy feedback to reduce frequency disagreement and control deviations from compatible occupancy targets, while respecting finite buffer capacity and oscillator-correction limits. This formulation makes transient buffer excursion, convergence and control activity joint criteria for evaluating a controller.

Existing bittide research establishes important foundations for this problem. PI performance analysis~\cite{lall2022resistance} relates closed-loop behaviour to controller gains and network resistance distances. Fluid occupancy analysis~\cite{lall2024bufferOccupancy} characterizes steady buffer levels, and reframing methods~\cite{lall2023bufferCentering} address buffer centering. Our earlier study, \emph{Stabilizing the bittide by Tuning Controllers Under Topology and Latency Constraints}, examined PI gain selection under topology and fixed-delay constraints. Building on these results, the question addressed here is how a local nonlinear Fuzzy PI law can reshape transient resource use under comparable scales, memory and update conventions. In particular, the connection between the nonlinear rule surface, a provable buffer-excursion benefit and behaviour under delivery perturbations requires explicit analysis and controlled comparison.

We propose a decentralized occupancy-based Fuzzy PI controller for bittide. Each node computes an oscillator-correction increment from its receiving-buffer error and the update-to-update change in that error. In the central input region studied here, the proposed surface strengthens its calibrated linear counterpart for same-sign inputs and attenuates its magnitude for opposing-sign inputs. This direction-dependent action provides flexibility at fixed calibrated PI gain coordinates. It motivates investigating whether the nonlinear correction can reduce transient buffer demands while retaining convergence, and determining how its effects on frequency disagreement and control activity depend on the operating conditions.

This paper presents the complete controller, establishes its con\-tinuous-clock realization and combines nonlinear analysis with evaluation to investigate that question. We derive the exact central-region map and its scale correspondence with native PI, and verify implementation equivalence when the nonlinear map is replaced by its calibrated linear counterpart. For a synchronous, zero-delay fluid reference, we establish an invariant input region and a common Lyapunov certificate. An outward interval enclosure of the nonlinear trajectory, followed by a contracting infinite-tail bound, provides an all-time buffer-excursion guarantee. For one prescribed three-node initialization at matched gain coordinates, the Fuzzy excursion is certified below 17.315 fluid frames, whereas PI attains an exact peak of approximately 23.403701 frames. This establishes a strict buffer-headroom benefit for that reference and initialization.

We separately evaluate the proposed law in continuous-phase Callisto,
including integer buffer events, local-clock polling and delayed
correction application. Matched-gain startup tests show smaller buffer
excursions for Fuzzy at three prescribed initial frequency mismatches.
We then give PI and fixed-shape Fuzzy the same 24-candidate,
four-case tuning budget, freeze the selected controllers, and execute
186 untouched tests using shared delivery profiles. In the two
original bounded-random regimes, the random-priority Fuzzy controller
reduces mean paired disagreement-RMS ratios by 16.78\% and 7.85\%
relative to the independently selected PI, with lower mean command
activity. The burst-priority Fuzzy controller also reduces IAE by
approximately 55.8\% in the larger held-out burst. These gains have
a defined operating scope: common-frequency shift is a cost, and a
smaller burst and a separately registered slower-random regime do not
share the improvement. The combined evidence establishes conditional
transient and synchronization benefits while keeping the reference
guarantee distinct from measured delayed execution.

The contributions of this paper are:
\begin{enumerate}
  \item A decentralized incremental Fuzzy PI controller for bittide,
  with complete inference rules, continuous-clock realization,
  implementation-checked PI calibration and an exact central-region
  nonlinear characterization.
  \item Invariant-domain, Lyapunov and transient-headroom analysis for
  the stated fluid reference, including positive-frequency and
  correction constraints and a strict all-time buffer-excursion
  benefit for a prescribed matched-gain initialization.
  \item Continuous-phase Callisto evaluation combining matched-gain
  startup stress, equal-budget independent tuning and untouched
  paired tests, identifying operating conditions with lower
  disagreement RMS or transient buffer demand and their associated
  performance tradeoffs.
\end{enumerate}

% The contributions of this paper are:

% \begin{enumerate}
%   \item \textbf{Fuzzy PI control for bittide.} A proposed decentralized occupancy-feedback controller with a fully specified inference law, initialization and continuous-clock actuation, together with PI calibration and exact nonlinear analysis that support controlled comparison.
%   \item \textbf{Convergence and transient buffer margins.} An invariant-domain and Lyapunov analysis for the specified fluid reference, with verified correction and positive-frequency bounds, and an all-time comparison proving a strict headroom benefit for a prescribed initialization at matched gain coordinates.
%   \item \textbf{Evaluation under delivery perturbations.} A continuous-clock Callisto study using shared traces, clean counterfactuals, complete-event buffer measurements and independent paired random seeds to characterize frequency, buffer-margin and control-activity tradeoffs.
% \end{enumerate}

% Historical fixed-preset introduction paragraph, retained for review.
% We separately evaluate the controllers in the continuous-phase Callisto model, with quantized buffer events, local-clock polling and control-execution delay. Controller updates change the slope of continuous phase trajectories; they do not replace the plant with a discrete firing-count model. Shared frame-delivery stress profiles include a latency burst after nominal centering and bounded random-latency traces. Over the declared measurement window beginning at burst onset, Fuzzy PI has lower integrated frequency-disagreement amplitude and command total variation, alongside a higher peak buffer excursion. In both random-latency regimes, the matched moderate PI baseline achieves lower disagreement RMS over the declared disturbance windows. These experiments characterize the performance tradeoffs accompanying nonlinear response shaping and distinguish the reference-model guarantee from its practical applicability under the declared execution assumptions.

% Expanded Background, 2026-11-03. Source-only revision; no PDF build.
\section{Background}
\label{sec:background}

\subsection{Logical Synchrony and the bittide Mechanism}

Logical synchrony describes relationships between distributed events without requiring a global physical clock. In bittide, neighbouring nodes exchange ordered frames, and elastic receiving buffers accommodate differences in their physical progress. The association between sending and receiving events provides the basis for a logical execution whose dependencies can be reasoned about independently of exact wall-clock timing. Equal physical phases are therefore not required for logical synchrony. Frequency agreement alone, however, does not establish the required event relationships: frame ordering, initialization and continued availability of the corresponding frames also matter. These distinctions follow the logical-synchrony and multiclock-network framework developed by Lall et al.~\cite{lall2024logicalSynchrony}.

The control problem supports this realization by regulating physical progress while retaining the prescribed frame relationships. Each node observes its own receiving buffers and adjusts its oscillator. Buffers consequently provide both the local feedback signal and finite storage for a temporary imbalance between delivery and consumption.

\subsection{Continuous Clocks and Receiving-Buffer Dynamics}

We consider a network of $n$ nodes. Node $i$ has continuous oscillator phase $\theta_i(t)$, measured in local cycles, and physical frequency

\begin{equation}
\dot\theta_i(t)=\omega_i(t)=\omega_i^u+u_i(t),
\end{equation}

where $\omega_i^u$ is its constant intrinsic frequency and $u_i(t)$ is the applied frequency correction. The correction can be signed; admissible physical operation requires $\omega_i(t)>0$. In the current Callisto model, corrections are held between local actuation events. Phase remains continuous when a correction changes, although its slope can change. Local sampling and integer frame events are embedded in this continuous-time plant.

For a receiving edge $e=(j\to i)$, let $\beta_e(t)$ be buffer occupancy in frames and $C_e$ its capacity. During normal operation without progress blocking,

\begin{align}
\beta_e(t)&=\beta_e(0)+A_{ji}^{\mathrm{arr}}(t)-N_i(t),\nonumber\\
N_i(t)&=\lfloor\theta_i(t)\rfloor-\lfloor\theta_i(0)\rfloor.
\label{eq:bg_count_balance}
\end{align}

Here $A_{ji}^{\mathrm{arr}}(t)$ counts arrivals during $(0,t]$, including frames already in flight at initialization, and $N_i(t)$ counts local consumption. This conservation relation explains why both oscillator mismatch and delivery changes affect occupancy. Section~\ref{sec:fpi_controller} specifies the local schedules, initialization and ordered arrival model used in our realization.

Fluid models replace integer frame counts by continuous phase increments. With constant physical path delay $\tau_{ji}$, the corresponding occupancy satisfies, between command discontinuities,

\begin{equation}
\dot{\widetilde\beta}_{ji}(t)
=\omega_j(t-\tau_{ji})-\omega_i(t),
\end{equation}

where $\widetilde\beta_{ji}$ denotes fluid occupancy. Thus occupancy integrates the imbalance between delayed source progress and local consumption. The fluid approximation facilitates analysis; its assumptions must be distinguished from those of integer-frame execution. This separation is also central to existing PI performance analysis~\cite{lall2022resistance} and buffer-occupancy modeling~\cite{lall2024bufferOccupancy}.

\subsection{Safety, Syntony and Buffer Centering}

Safety is a property of the complete trajectory. For the normal operating regime studied here, we seek

\begin{equation}
0<\beta_e(t)<C_e,\qquad
\omega_i(t)>0,\qquad |u_i(t)|\le u_{\max},
\end{equation}

for every edge, node and time, where $u_{\max}>0$ is the correction limit. The available lower and upper buffer reserves are $\beta_e(t)$ and $C_e-\beta_e(t)$, respectively. A convergent trajectory can still exhaust either reserve before reaching its limit. Avoiding such exhaustion through control is distinct from invoking a protection mechanism that blocks logical progress.

Define the physical frequency spread as

\begin{equation}
\Delta_\omega(t)=\max_i\omega_i(t)-\min_i\omega_i(t).
\end{equation}

For nominal fluid operation with fixed conditions, we seek $\Delta_\omega(t)\to0$, together with $\beta_e(t)\to\beta_e^*$ for compatible occupancy targets $\beta_e^*$. Frequency agreement is the nominal syntonization objective; it does not imply zero phase differences or tracking an externally prescribed frequency. Without an external reference, the common limiting frequency $\omega_\infty$ is determined by the closed loop. Likewise, arbitrary per-edge targets cannot be assumed compatible with a fixed marking, meaning the frame-count offsets determined by initialization, and with the selected initial state.

If a transient frequency-spread tolerance $\Delta_{\mathrm{safe}}>0$ is specified, then $\Delta_\omega(t)<\Delta_{\mathrm{safe}}$ is an additional engineering constraint. It is not the definition of logical synchrony. Under persistent delivery perturbations, we assess bounded deviations, residual disagreement and buffer reserves rather than assume exact disturbance-free centering.

\subsection{Delivery Perturbations and Local Feedback}

A latency increase affects the delivery process before a receiver necessarily observes it through its next control measurement. If arrivals pause while consumption continues, the conservation relation shows that the receiving buffer decreases; subsequent release of accumulated frames can instead increase occupancy. The disturbance therefore acts through frame availability and accumulated count imbalance. A controller cannot recreate an undelivered frame, and its response is constrained by measurement and actuation timing. Our experiments use declared order-preserving delivery traces so that controller comparisons share the same external stress definition. The exact arrival construction is given in Section~\ref{sec:fpi_controller}.

\subsection{Existing Control Foundations and the Present Study}

Existing work analyzes occupancy-feedback PI performance through controller gains and graph resistance distances, derives limiting fluid occupancies, and develops reframing methods that separate frequency regulation from occupancy positioning. These are complementary foundations: PI performance~\cite{lall2022resistance}, steady buffer occupancy~\cite{lall2024bufferOccupancy}, and buffer centering~\cite{lall2023bufferCentering}. Their conclusions apply under their respective modeling assumptions; in particular, continuous-time gain conditions cannot be transferred unchanged to sampled execution with delay.

This paper builds on those foundations by proposing a local incremental Fuzzy PI law and analyzing how its nonlinear increments reshape the closed-loop response at explicitly calibrated PI scales. Section~\ref{sec:fpi_controller} defines the complete controller and establishes that calibration. Section~\ref{sec:btd_reference} derives sufficient safety and convergence conditions for a declared synchronous, zero-delay fluid reference, followed by an all-time buffer-excursion comparison for one prescribed initialization. Evaluation in continuous-phase Callisto separately examines delivery perturbations under the actual local update and frame-count conventions.

% Expanded, input-ready Section III. No standalone preamble or PDF build.
\section{Proposed Fuzzy PI Control for bittide and Nonlinear Analysis}
\label{sec:fpi_controller}
We propose a decentralized, occupancy-based incremental Fuzzy PI controller
for bittide. Each node combines receiving-buffer error with its
update-to-update change to compute an oscillator-correction increment.
This section specifies the complete controller and its continuous-clock
interface, derives an implementation-consistent PI calibration, and
establishes the nonlinear properties used in
Section~\ref{sec:btd_reference}. The physical plant, integer frame counts
and sampled controller state are kept distinct throughout the derivation.

\subsection{Continuous-clock realization of the proposed controller}

\label{subsec:fpi_laws}
The plant is the current continuous-phase Callisto model within the bittide
framework~\cite{lall2024logicalSynchrony}. For constant
intrinsic oscillator frequency $\omega_i^u$ and applied correction $u_i(t)$,
\begin{align}
 \dot\theta_i(t)&=\omega_i^u+u_i(t),\nonumber\\
 u_i(t)&=u_{i,k},\qquad t\in[t^a_{i,k},t^a_{i,k+1}),\nonumber\\
 \theta_i(t)&=\theta_i(t^a_{i,k})+
       (\omega_i^u+u_{i,k})(t-t^a_{i,k}).
 \label{eq:fpi_continuous_plant}
\end{align}
The superscript $a$ denotes actuation, and $k$ indexes that node's updates.
Normal executions require $\omega_i^u+u_i(t)>0$, so the phase has a
monotone inverse. Physical phase is continuous at an update; its slope can change. Requested
correction $q_{i,k}$ is computed from the measured buffers and applied as
$u_{i,k}=\operatorname{sat}_{u_{\max}}(q_{i,k})$, with
$\operatorname{sat}_b(z)=\min\{b,\max\{-b,z\}\}$ and $u_{\max}>0$.
The proposed controller has jump updates inside a continuous physical plant;
it is not implemented as an ordinary differential equation for its accumulator.

Let $p>0$ be the polling period and $d_{\rm ctrl}\ge0$ the execution delay,
both measured in local-clock phase units. Measurement time $t^m_{i,k}$
and actuation time $t^a_{i,k}$ satisfy
\begin{align}
 \theta_i(t^m_{i,k})&=\theta_i(0)+k\,p,\nonumber\\
 \theta_i(t^a_{i,k})&=\theta_i(0)+k\,p+d_{\rm ctrl},\nonumber\\
 \int_{t^m_{i,k}}^{t^a_{i,k}}(\omega_i^u+u_i(s))\,ds&=d_{\rm ctrl}.
 \label{eq:fpi_local_schedule}
\end{align}
The product $k\,p$ in these scheduling equations means the integer index
$k$ multiplied by the poll period $p$; it is not the proportional gain
$k_p$. With zero initial correction, $t^m_{i,0}=0$ and
$t^a_{i,0}=d_{\rm ctrl}/\omega_i^u$. Subsequent actuation intervals equal
$p/(\omega_i^u+u_{i,k})$. The delay is not automatically a constant
physical time, and its age must not be computed using the new command
that is applied only at the end of that age interval.

In normal operation, each integer source-phase crossing sends a frame.
With the progress guards open, frame conservation gives
\begin{equation}
 \beta_{j\to i}(t)=\beta_{j\to i}(0)+A^{\rm arr}_{ji}(t)-N_i(t),
 \label{eq:fpi_count_balance}
\end{equation}
where $A^{\rm arr}_{ji}(t)$ counts arrivals during $(0,t]$, including
initial in-flight frames, and
$N_i(t)=\lfloor\theta_i(t)\rfloor-\lfloor\theta_i(0)\rfloor$ counts
local consumption. For a constant physical delay $\tau_{ji}\ge0$,
\begin{equation}
 A^{\rm arr}_{ji}(t)=\lfloor\theta_j(t-\tau_{ji})\rfloor
                  -\lfloor\theta_j(-\tau_{ji})\rfloor.
 \label{eq:fpi_fixed_delay_arrivals}
\end{equation}
Thus arrival counts incorporate the declared negative-time history.

For constant physical path delay $\tau_{ji}$, receiving occupancy is
\begin{equation}
 \beta_{j\to i}(t)=\lambda^{\rm UGN}_{ji}
 +\lfloor\theta_j(t-\tau_{ji})\rfloor-\lfloor\theta_i(t)\rfloor.
 \label{eq:fpd_continuous_fifo}
\end{equation}
Here $\lambda^{\rm UGN}_{ji}$ is the fixed integer offset in the absolute
phase-count convention. It is determined by occupancy and prehistory:
\begin{equation}
 \lambda^{\rm UGN}_{ji}=\beta_{j\to i}(0)
 -\lfloor\theta_j(-\tau_{ji})\rfloor+\lfloor\theta_i(0)\rfloor.
 \label{eq:fpd_ugn_init}
\end{equation}
The relative-count logical marking is
$\lambda^{\rm LSN}_{ji}=\lambda^{\rm UGN}_{ji}
+\lfloor\theta_j(0)\rfloor-\lfloor\theta_i(0)\rfloor$;
the two coincide when the initial phase counts are equal.

For changing delivery times, let $s_{j,m}=\theta_j^{-1}(m)$ be the send
time of source tick $m$, and $a_{ji}(s)$ its causal ordered arrival map.
For the desired send-time delay profile $\tau_{ji}(s)$, the declared
running-maximum envelope is
\[
 a_{ji}(s)=\sup_{r\le s}\{r+\tau_{ji}(r)\},
\]
including the constant-delay negative-time history.
Define the delivered tick frontier
\begin{equation}
 C^{\rm arr}_{ji}(t)=\max\{m\in\mathbb Z:a_{ji}(s_{j,m})\le t\},
 \quad\beta_{j\to i}(t)=\lambda^{\rm UGN}_{ji}
       +C^{\rm arr}_{ji}(t)-\lfloor\theta_i(t)\rfloor.
 \label{eq:fpd_ordered_arrivals}
\end{equation}
The negative-time history specifies initial in-flight frames. The declared
running-maximum arrival envelope preserves order and can release a batch
after a desired-delay decrease; it imposes no finite serialization rate.
At an upward-delay discontinuity, a supremum of continuous source times
need not itself have arrived. Counting delivered tick indices avoids
silently including such a boundary frame.

If a protection guard blocks logical progress, the runtime can retain a
separate continuous data phase $\psi_i$ with
$\dot\psi_i=g_i(t)\dot\theta_i$, where $g_i\in\{0,1\}$ is the progress
gate. Buffer counts then use $\psi_i$, whereas scheduling still uses
$\theta_i$. In the ungated normal regime, with matching initial phases, data phase and physical phase agree. Guard operation does not stop the oscillator.

Only for the reference proof do we replace integer occupancy by a
fluid counterpart. At constant delay, define
$\widetilde\beta_{ji}(t)=\widetilde\lambda_{ji}
+\theta_j(t-\tau_{ji})-\theta_i(t)$, where
$\widetilde\lambda_{ji}$ is an additive offset chosen for the specified
fluid initialization. Between command discontinuities,
\begin{equation}
 \dot{\widetilde\beta}_{ji}(t)
 =\omega_j^u+u_j(t-\tau_{ji})-\omega_i^u-u_i(t).
 \label{eq:fpi_fluid_counterpart}
\end{equation}
This derivative belongs to the fluid counterpart, not the actual integer
occupancy. Replacing floor counts by phases does not automatically preserve
an arbitrary initialization; the compatible zero-delay marking used in
Section~\ref{sec:btd_reference} is specified there.

\subsection{Notation and the measured error}

Node $i$ has incoming-neighbour set $\mathcal N_i^-$ and $d_i=|\mathcal N_i^-|>0$ incoming channels. All incoming sums below range over $j\in\mathcal N_i^-$. This section assumes fixed active
incoming sets: the native link-status mask is one on all configured inputs.
The Fuzzy percentage wrapper averages its configured inputs, so changing
link masks would require an explicit corresponding measurement policy.
At the $k$th local measurement,
$\beta_{ji,k}$ is the measured occupancy of channel $j\to i$, and
$\beta^*_{ji}$ is its fixed target. Use
\begin{equation}
 E_{i,k}=\sum_{j\in\mathcal N_i^-}(\beta_{ji,k}-\beta^*_{ji}).
 \label{eq:fpi_error}
\end{equation}
The sign is actual minus target. If incoming buffers are too full, a
positive frequency correction increases local consumption. This explains
the positive PI gains with this error convention. The physical frequency is $\omega_i=\omega_i^u+u_i$, where $\omega_i^u$ is the constant intrinsic frequency; $u_i$ is a correction, not the entire frequency.

\begin{center}
\begin{tabular}{@{}ll@{}}
\toprule
Symbol & Meaning \\
\midrule
$E_{i,k}$ & Sum error in frames \\
$e_{i,k}$ & Target-minus-measured error in percentage points \\
$I_{i,k}$ & Native PI sum of sampled frame errors \\
$A_{i,k}$ & Fuzzy accumulator: total requested correction \\
$q_{i,k},u_{i,k}$ & Requested and applied frequency corrections \\
$h$ & Fixed native PI integration scale \\
$X_{i,k},Y_{i,k}$ & Scaled Fuzzy error and sample-difference inputs \\
$F_0,\mathcal R_F$ & Linear counterpart and residual $F-F_0$ \\
\bottomrule
\end{tabular}
\end{center}

The event index $k$ is local to a node. The normalization factor $\gamma_i$ converts a frame-sum error into percentage points. The maximum correction $u_{\max}>0$ is the actuator limit; $q$ and $u$ remain distinct if clipping acts. Two nodes need not measure at the
same physical times. The subscript $ji$ abbreviates edge $j\to i$.
Suppressing $i$ below is a notational convenience, not a synchrony
assumption; in that notation $\gamma=\gamma_i$.

\subsubsection{Deriving the percentage-to-token relation}
For common incoming capacity $C$, define
\[
 \bar p_{i,k}=\frac1{d_i}\sum_j\frac{100\beta_{ji,k}}{C},\qquad
 r_i=\frac1{d_i}\sum_j\frac{100\beta^*_{ji}}{C}.
\]
Then subtraction and factoring give
\begin{align}
 e_{i,k}&=r_i-\bar p_{i,k}\\
 &=-\frac{100}{d_iC}\sum_j(\beta_{ji,k}-\beta^*_{ji})
 =-\gamma_i E_{i,k},\qquad \gamma_i=\frac{100}{d_iC}.
 \label{eq:fpi_percent}
\end{align}
For midpoint targets this is the original absolute-percentage input with
$r_i=50$. For other targets the runtime can supply a centered percentage
$r_i+d_i^{-1}\sum_j100(\beta_{ji,k}-\beta^*_{ji})/C$; subtracting it from
$r_i$ yields the same $e_{i,k}$. A mismatched setpoint instead adds a
constant bias and invalidates the scalar relation used for calibration.

With unequal capacities and a consistently matched target reference, the
percentage average gives
\[
 e_{i,k}=-\frac{100}{d_i}\sum_j
 \frac{\beta_{ji,k}-\beta^*_{ji}}{C_{ji}}.
\]
This is a weighted error and generally cannot be written as one normalization constant
times the unweighted native sum. Matching such a controller requires a
consistent measurement weighting or a different baseline.

\subsection{Native PI: update order, units and the first call}

The occupancy-feedback PI baseline is grounded in the bittide analysis
of~\cite{lall2022resistance}. The executed native update is an
integrate-then-output recursion:
\begin{equation}
 I_k=I_{k-1}+E_k,\qquad
 q_k^{\rm PI}=k_pE_k+\kappa_I I_k,\qquad \kappa_I=k_i h.
 \label{eq:fpi_native}
\end{equation}
The sum $I_k$ contains no physical time multiplier; $h$ is incorporated in
$\kappa_I$. In the present runtime,
\begin{equation}
 h=\frac{p}{\omega_{\rm base}}.
 \label{eq:fpd_h}
\end{equation}
Here $p$ is the local polling period and $\omega_{\rm base}$ is the fixed implementation reference frequency. The scale $h$ stays fixed during the configured run. Local polling intervals measured
in physical time can vary as the oscillator changes; replacing $h$ by
those intervals would implement a different PI law.

Subtract the previous requested output:
\begin{align}
 q_k^{\rm PI}-q_{k-1}^{\rm PI}
 &=k_p(E_k-E_{k-1})+\kappa_I(I_k-I_{k-1})\\
 &=k_p(E_k-E_{k-1})+\kappa_I E_k.
 \label{eq:fpi_pi_increment}
\end{align}
This is the correct increment to reproduce. It contains the present
$E_k$ in the integral term because integration precedes output.

At zero initial sum, the first call returns
\begin{equation}
 q_0^{\rm PI}=(k_p+\kappa_I)E_0.
 \label{eq:fpd_first}
\end{equation}
No physical prehistory measurement is needed to express this algebra:
one may use the fictitious previous error $E_{-1}=0$ and previous requested
output $q_{-1}=0$. That bookkeeping choice is separate from deriving a
plant identity such as $x_k-x_{k-1}=H w_{k-1}$, which requires a genuine
elapsed plant interval.

\subsection{The proposed Fuzzy PI controller}

The proposed local incremental control law combines the present error
and its sample difference as follows:
\begin{align}
 X_k&=K_e e_k,&Y_k&=K_{\Delta e}(e_k-e_{k-1}),\\
 A_k&=A_{k-1}+K_uF(X_k,Y_k),&q_k^{\rm F}&=A_k.
 \label{eq:fpi_fuzzy_law}
\end{align}
The gain signs used here are $K_e,K_{\Delta e}>0$ and $K_u<0$. Thus high occupancy
produces negative $e$, and a negative surface output produces a positive
correction increment. The quantity $e_k-e_{k-1}$ is not divided by elapsed
time; calling it a time derivative would obscure its sampling dependence.

\subsubsection{Memberships and rule weights}
The labels NL, NS, ZE, PS and PL mean negative large, negative small, zero, positive small and positive large, respectively. Each membership $\varphi_a(z)\in[0,1]$ is an activation degree; a singleton is a numerical rule consequent. Let $[z]_+=\max\{z,0\}$ and $\operatorname{clip}_{[0,1]}z$
denote projection onto $[0,1]$. The exact memberships can be written
\begin{align}
 \varphi_{\rm NL}(z)&=\operatorname{clip}_{[0,1]}\frac{-z-5}{5},&
 \varphi_{\rm NS}(z)&=[1-|z+5|/5]_+,\nonumber\\
 \varphi_{\rm ZE}(z)&=[1-|z|/5]_+,&
 \varphi_{\rm PS}(z)&=[1-|z-5|/5]_+,\nonumber\\
 \varphi_{\rm PL}(z)&=\operatorname{clip}_{[0,1]}\frac{z-5}{5}.
 \label{eq:fpi_full_memberships}
\end{align}
Table~\ref{tab:fpi_full_rules} gives all 25 consequent rules.
\begin{table}[t]
\caption{Rule consequents of the proposed bittide Fuzzy PI controller.}
\label{tab:fpi_full_rules}
\centering
\begin{tabular}{@{}c|ccccc@{}}
\toprule
$X\backslash Y$ & NL & NS & ZE & PS & PL\\
\midrule
NL & NL & NL & NL & NS & ZE\\
NS & NL & NL & NS & ZE & PS\\
ZE & NL & NS & ZE & PS & PL\\
PS & NS & ZE & PS & PL & PL\\
PL & ZE & PS & PL & PL & PL\\
\bottomrule
\end{tabular}
\end{table}
Order the labels as NL, NS, ZE, PS and PL, with label indices
$a,b\in\{1,\ldots,5\}$. With singleton values
$(s_1,\ldots,s_5)=(-5,-2.5,0,2.5,5)$ and
$\ell(a,b)=\min\{5,\max\{1,a+b-3\}\}$, compute
\begin{equation}
 w_{ab}=\min\{\varphi_a(X),\varphi_b(Y)\},\qquad
 F(X,Y)=\frac{\sum_{a,b}w_{ab}s_{\ell(a,b)}}{\sum_{a,b}w_{ab}}.
 \label{eq:fpi_rules}
\end{equation}
The membership partition gives at least one positive input membership on
each axis, hence a positive denominator at finite inputs. Individual rules
are summed even when they share a consequent. Replacing this accumulation
by maximum aggregation changes the map and its derivation.

The weighted-average property gives $|F|\le5$ and therefore
$|A_k-A_{k-1}|\le5|K_u|$. This is a per-update bound, not a bound on $A_k$
over an arbitrarily long execution. The raw Fuzzy module itself has no
final correction clipping.
Initialize $A_{i,-1}=0$ and $e_{i,-1}=0$. At each measurement, form
$e_{i,k}$ and $(X_{i,k},Y_{i,k})$, evaluate the rule weights and normalized
surface, increment $A_{i,k}$, and store $e_{i,k}$ as the previous error.
The requested correction is then limited and applied at $t^a_{i,k}$, and
held until the next actuation. In the normal-run policy considered here,
external clipping retains the raw accumulator and anti-windup is disabled.
The controller does not overwrite $A_{i,k}$ with the applied command.
The membership functions, rule consequents and gains stay fixed during a
run; the correction changes with the measured input pair. The law requires
neither an observer nor an FFP mode.


\subsection{Exact PI Calibration, Including Initialization}
\label{subsec:fpi_calibration}

Define the globally linear counterpart
\begin{equation}
 F_0(X,Y)=\frac{X+Y}{2}.
 \label{eq:fpi_f0}
\end{equation}
Replacing $F$ by this map retains the same controller accumulator.
Using $e=-\gamma E$ gives
\begin{align}
 A^0_k-A^0_{k-1}
 &=\frac{K_u}{2}\{K_e e_k+K_{\Delta e}(e_k-e_{k-1})\}\\
 &=-\frac{K_u\gamma K_e}{2}E_k
   -\frac{K_u\gamma K_{\Delta e}}2(E_k-E_{k-1}).
 \label{eq:fpi_substitution}
\end{align}
Matching the two independent coefficients in~\eqref{eq:fpi_pi_increment}
yields
\begin{equation}
 k_p=-\frac{K_u\gamma K_{\Delta e}}2,\qquad
 \kappa_I=-\frac{K_u\gamma K_e}2.
 \label{eq:fpi_matching}
\end{equation}
Solving for two Fuzzy gains with chosen $K_e>0$ gives
\begin{equation}
 K_u=-\frac{2k_ih}{\gamma K_e},\qquad
 K_{\Delta e}=\frac{K_ek_p}{k_ih}.
 \label{eq:fpi_design_gains}
\end{equation}
The formulas assume $k_i h>0$. Pure proportional control is a separate
limiting case and should not be inserted into this inverse mapping with
a zero denominator.

\begin{btdproposition}[Event-level linear-counterpart equivalence]
\label{prop:fpi_equivalence}
For any measurement sequence, suppose~\eqref{eq:fpi_matching} holds and
initialize
\begin{equation}
 A^0_{-1}=k_pE_{-1}+\kappa_I I_{-1},\qquad e_{-1}=-\gamma E_{-1}.
 \label{eq:fpi_initialization}
\end{equation}
Then
$A^0_k=k_pE_k+\kappa_I I_k$ for every $k\ge0$.
\end{btdproposition}
\begin{proof}
Assume the relation holds before update $k$. Substitution gives
\begin{align*}
 A^0_k
 &=k_pE_{k-1}+\kappa_I I_{k-1}
   +k_p(E_k-E_{k-1})+\kappa_I E_k\\
 &=k_pE_k+\kappa_I(I_{k-1}+E_k)
 =k_pE_k+\kappa_I I_k.
\end{align*}
The stated initialization supplies the induction base, including $k=0$.
\end{proof}

\subsubsection{Why zero constructors work for a nonzero first error}
Take $I_{-1}=A^0_{-1}=e_{-1}=0$, corresponding to $E_{-1}=0$.
The first Fuzzy-reference input is
\[
 (X_0,Y_0)=(-\gamma K_eE_0,-\gamma K_{\Delta e}E_0),
\]
so
$A^0_0=(k_p+\kappa_I)E_0$, exactly as in~\eqref{eq:fpd_first}.
If only Fuzzy is warmed up to $e_{-1}=e_0$ while $A^0_{-1}=0$, its
difference input vanishes and the output is merely $\kappa_I E_0$.
The missing $k_pE_0$ term can create a large apparent controller difference
unrelated to fuzzy inference.

\subsubsection{A numerical calibration from the present configuration}
For $d=2$, $C=2048$, $p=1000$, $\omega_{\rm base}=29999/30000$,
$k_p=2\times10^{-6}$, $k_i=5\times10^{-11}$ and $K_e=0.1$,
\begin{align}
 \gamma&=\frac{25}{1024},&h&=\frac{30000000}{29999},\nonumber\\
 \kappa_I&=\frac3{59998000},&
 K_u&=-\frac{768}{18749375},\nonumber\\
 K_{\Delta e}&=\frac{29999}{7500}
 \simeq3.9998666667.
 \label{eq:fpd_example_gains}
\end{align}
For a one-frame first error, the matched PI and $F_0$ corrections are both
$2.0500016667\times10^{-6}$. Warming up only the Fuzzy previous error
instead gives approximately $5.0001666722\times10^{-8}$. The example
illustrates an initialization mismatch; it is not a new plant simulation.

\subsubsection{Execution-layer conditions}
Identical raw outputs give identical applied outputs under a common
limiter when raw states are retained and anti-windup is disabled.
Tracking or reset requires compatible state changes. For example, after
processing $E_k$, setting
\[
 A^0_k=u_{\rm held},\qquad
 I_k=\frac{u_{\rm held}-k_pE_k}{\kappa_I},\qquad
 e_k=-\gamma E_k
\]
restores the state relation for the next ordinary update. A branch designed
to produce a held output on an immediate extra update requires its own
integrate-first adjustment; the above is not an instruction to add an
extra call. Limiting the internal reference $F_0$ to $[-5,5]$ also changes
the controller, even if the external correction limiter is shared.

Under a deterministic causal plant, the same initial history, measured
inputs, actuation order, external trace and compatible resets imply
closed-loop equality by induction over events. No globally synchronous
physical grid is required for this algebraic equivalence. By contrast,
matching only labels such as ``same gains'' or ``same random seed'' does
not establish these conditions.

\subsection{Deriving the Exact Central Nonlinear Surface}
\label{subsec:fpi_mechanism}

For $0<r\le5$, define
\begin{align}
 \Omega_r&=\{(X,Y):|X|+|Y|\le r\},\nonumber\\
 m&=\min(|X|,|Y|),\qquad s=|X|+|Y|.
 \label{eq:fpi_domain}
\end{align}
The same-sign interaction is
\begin{equation}
 J(X,Y)=\begin{cases}\operatorname{sgn}(X)m,&XY>0,\\0,&XY\le0.\end{cases}
 \label{eq:fpi_interaction}
\end{equation}
For points in $\Omega_5$, put $a=|X|/5$, $b=|Y|/5$.
Because $a+b\le1$, each input has only its central and sign-adjacent
membership active. Their four joint min weights are
\begin{align}
 w_{00}&=\min(1-a,1-b)=1-\max(a,b),\nonumber\\
 w_{s0}&=\min(a,1-b)=a,\nonumber\\
 w_{0s}&=\min(1-a,b)=b,\nonumber\\
 w_{ss}&=\min(a,b).
 \label{eq:fpd_four_weights}
\end{align}
The subscripts $s$ denote the input's active sign label, not the scalar
$s=|X|+|Y|$. Zero-input boundaries follow by continuity.

The denominator is therefore
\begin{equation}
 D=1-\max(a,b)+a+b+\min(a,b)
  =1+2\min(a,b)=1+2m/5.
 \label{eq:fpd_denominator}
\end{equation}
The central-central consequent is zero. The sign-central and
central-sign rules contribute $(5/2)\operatorname{sgn}(X)a=X/2$
and $(5/2)\operatorname{sgn}(Y)b=Y/2$.
For equal signs, the joint consequent is
$5\operatorname{sgn}(X)$ and contributes $\operatorname{sgn}(X)m$;
for opposing signs it is zero. The four quadrants therefore give
\begin{equation}
 F(X,Y)=\begin{cases}
 \dfrac{F_0+m}{1+2m/5},&X>0,\ Y>0,\\[4pt]
 \dfrac{F_0-m}{1+2m/5},&X<0,\ Y<0,\\[4pt]
 \dfrac{F_0}{1+2m/5},&X>0,\ Y<0,\\[4pt]
 \dfrac{F_0}{1+2m/5},&X<0,\ Y>0.
 \end{cases}
 \label{eq:fpi_quadrants}
\end{equation}
On either axis $m=0$, giving $F=F_0$. These cases establish the following
identity, including boundaries by continuity.
\begin{btdproposition}[Exact central surface]
\label{prop:fpi_central}
For every $(X,Y)\in\Omega_5$,

\begin{equation}
 F(X,Y)=\frac{(X+Y)/2+J(X,Y)}{1+2m/5},\qquad
 J(X,Y)=\begin{cases}\operatorname{sgn}(X)m,&XY>0,\\0,&XY\le0.
 \end{cases}
 \label{eq:fpi_central}
\end{equation}
\end{btdproposition}
This is an exact rule-level identity on $\Omega_5$, not a fitted
approximation. Outside this diamond, the cross weights in
\eqref{eq:fpd_four_weights} need not equal $a$ and $b$, so the formula
must not be extended without a new case analysis.

\subsection{Sign-conditioned response, sector and nonlinear geometry}

\subsubsection{Same-sign enhancement}
When $XY>0$, $F_0$ has magnitude $s/2$ and the same sign as $X$.
Equation~\eqref{eq:fpi_central} gives
\begin{equation}
 \frac{F}{F_0}=\frac{1+2m/s}{1+2m/5}.
 \label{eq:fpi_same_sign}
\end{equation}
Since $s\le5$, its numerator is at least the denominator. Since
$m\le s/2$, the numerator is at most 2. Thus $1\le F/F_0\le2$.
Define the normalized nonlinear residual by
$\mathcal R_F(X,Y)=F(X,Y)-F_0(X,Y)$. Subtracting $F_0$ gives
\begin{equation}
 \mathcal R_F=
 \frac{\operatorname{sgn}(X)m(1-s/5)}{1+2m/5},\qquad XY>0.
 \label{eq:fpd_same_residual}
\end{equation}
For instance, $F(0.1,0.1)=0.2/1.04\simeq0.19230769$, compared with
$F_0=0.1$. Both are surface values; the actual increment multiplies by
negative $K_u$.

\subsubsection{Opposing-sign attenuation and zero-output cases}
When $XY<0$, the interaction vanishes, so
\begin{equation}
 F=\frac{F_0}{1+2m/5},\qquad
 \mathcal R_F=-\frac{2m/5}{1+2m/5}F_0.
 \label{eq:fpi_opposite_sign}
\end{equation}
For $(X,Y)=(0.1,-0.05)$, $F_0=0.025$ and
$F=0.025/1.02\simeq0.024509804$. On either axis $m=0$ and $F=F_0$.
On the anti-diagonal $Y=-X$, both outputs are exactly zero; a ratio should
not be evaluated there.

\begin{btdcorollary}[Regional pointwise sector]
\label{cor:fpi_sector}
For $(X,Y)\in\Omega_r$, $0<r\le5$, one can write
\begin{equation}
 F=\alpha(X,Y)F_0,\qquad
 \underline\alpha_r=\frac1{1+r/5}\le\alpha(X,Y)\le2.
 \label{eq:fpi_sector}
\end{equation}
\end{btdcorollary}
\begin{proof}
Same-sign ratios lie in $[1,2]$. For opposing signs with $F_0\ne0$,
$m\le r/2$ gives $1/(1+r/5)\le F/F_0\le1$.
The axes agree with $F_0$. If $F_0=0$, both outputs vanish and any
admissible $\alpha$ preserves the equality.
\end{proof}
The equivalent division-free inequality is
\begin{equation}
 [F-\underline\alpha_r F_0][2F_0-F]\ge0.
 \label{eq:fpi_sector_product}
\end{equation}
In particular, $\Omega_1$ gives $\alpha\in[5/6,2]$, while
$\Omega_5$ gives $[1/2,2]$. These bounds describe pointwise values of
the specific nonlinear map. They are not a gain-scheduled stability proof;
the closed loop must retain the input region and accommodate different
node-dependent, time-varying $\alpha$ values.

\subsubsection{Why an axis match is not a Jacobian linearization}
The same-sign term $J$ is positively homogeneous. For $t>0$ and a fixed
direction $v=(v_1,v_2)$, the central formula yields
\begin{equation}
 F'(0;v)=\lim_{t\downarrow0}\frac{F(tv_1,tv_2)}t
 =\frac{v_1+v_2}{2}+J(v_1,v_2).
 \label{eq:fpd_direction}
\end{equation}
Consequently $F'(0;(1,0))=F'(0;(0,1))=1/2$, but
$F'(0;(1,1))=2$. A linear derivative with the two axis values would
instead give 1 on $(1,1)$, a contradiction. The map is nonlinear and not
Fr\'echet differentiable at the origin, despite exact agreement with
$F_0$ on the coordinate axes. ``Axis-calibrated linear counterpart'' is
therefore the appropriate term.

\subsubsection{Where the central interaction comes from}
As an analytic counterfactual, replace min rule weights by
$w^{\rm prod}_{ab}=\varphi_a(X)\varphi_b(Y)$, retaining the membership functions
and singletons. On $|X|,|Y|\le5$, the active consequents add their two
single-input values and the denominator is
$(\sum_a\varphi_a(X))(\sum_b\varphi_b(Y))=1$. Thus
\begin{equation}
 F_{\rm prod}(X,Y)=X/2+Y/2=F_0(X,Y).
 \label{eq:fpd_product}
\end{equation}
The central nonlinear interaction is tied to normalized min activation,
not merely to linguistic labels. The identity isolates the role of min activation; it does not assert
linearity beyond the central square or an experimental comparison of
product and min inference.

\subsection{Nonlinear memory and what it does not identify}

For identical scaled input pairs $(X_k,Y_k)$, including consistent
previous-error initialization, let $A_k^{\rm F}$ use $F$ and
$A_k^0$ use $F_0$, with equal initial accumulators. Subtracting their
updates gives
\[
 A_k^{\rm F}-A_k^0=(A_{k-1}^{\rm F}-A_{k-1}^0)
                  +K_u\mathcal R_F(X_k,Y_k).
\]
Telescoping yields
\begin{equation}
 A_k^{\rm F}-A_k^0=K_u\sum_{\ell=0}^{k}\mathcal R_F(X_\ell,Y_\ell).
 \label{eq:fpi_memory}
\end{equation}
With zero native/Fuzzy constructor states and the calibrated parameters,
the same-input bookkeeping can also be written
\begin{equation}
 A_k^{\rm F}=k_pE_k+\kappa_I\sum_{\ell=0}^{k}E_\ell
             +K_u\sum_{\ell=0}^{k}\mathcal R_F(X_\ell,Y_\ell).
 \label{eq:fpd_accounting}
\end{equation}
The first two terms form a PI-like shadow evaluated on the Fuzzy
trajectory's own inputs. The last term records accumulated nonlinear
correction. If later inputs lie on the axes, the new residual is zero
but the accumulated difference can remain nonzero.

This identity requires no accumulator projection/reset. Pure external
clipping with retained raw states leaves the raw identity intact, but the
applied difference must be evaluated through the limiter. Tracking/reset
requires explicit additional state-change terms. Moreover, independently
evolving PI and Fuzzy closed loops generate different measurements and
possibly different physical event schedules. A same-input shadow identity
therefore cannot be relabelled as their observed closed-loop difference.

\subsection{Interpretation and the bridge to the stability section}

The derivation provides implementation-consistent PI calibration,
the exact central nonlinear map, and a regional sector suitable for a
closed-loop certificate. Section~\ref{sec:btd_reference} uses these properties after stating its
plant, sampling, latency and initialization assumptions and proving that
the sector's input region remains invariant. The present algebra does not
change the continuous-clock Callisto plant into a discrete-time plant;
it describes the controller updates embedded in that execution.

Same-sign scaled inputs indicate alignment of the two input channels.
They do not always mean increasing absolute error: a transition from
$e_{k-1}=-10$ to $e_k=1$ has $e_k(e_k-e_{k-1})>0$ even though the
absolute error decreased. Likewise, opposing-sign attenuation concerns a
single increment relative to $F_0$, not an established network jitter
attenuation property. Frequency disagreement, buffer excursion, common
frequency and control activity must be evaluated separately.


% Expanded paper insert, 2026-11-03. No standalone document or bibliography.
% Source-text revision only; no PDF generation or simulation is required.
\section{Reference-Model Convergence and Buffer Margins}
\label{sec:btd_reference}

We connect the nonlinear controller of Section~\ref{sec:fpi_controller}
to four separate properties: invariance of its analyzed input region,
convergence of frequency and buffer disagreements, strict buffer-capacity
safety, and positive physical oscillator frequency. We then establish a
strict headroom comparison with its calibrated PI baseline for one
specified initialization. The proof concerns an explicitly defined
continuous fluid reference. Actual Callisto retains asynchronous local
sampling, delivery delays and integer frame events, which are not
silently removed from its experimental model.

\subsection{Reference assumptions and continuous-time model}
\label{subsec:btd_ref_model}

\subsubsection{What the reference retains and idealizes}
The reference retains the controller of Section~\ref{sec:fpi_controller}, continuous oscillator phases, local occupancy
feedback, the native integrate-before-output PI convention and the proposed
incremental normalized-min Fuzzy surface. It assumes zero path delay,
zero control-execution delay, fluid rather than integer occupancy and a
common physical sampling grid $t_k=kH$, $H>0$. A zero-order hold (ZOH)
means that the applied frequency correction is constant on each interval
$[t_k,t_{k+1})$. These are explicit analysis assumptions; local-clock
polling in actual Callisto is not asserted to have this common grid.
The equations sampled at $kH$ below are exact endpoint relations for that
continuous ZOH flow. They are neither numerical Euler integration nor a
discrete firing-count replacement of the plant.

Let $\mathcal N_i^-$ be the fixed incoming-neighbour set and
$d_i=|\mathcal N_i^-|$. The graph is connected, undirected and unit-weight,
with a receiving channel in both directions. Its symmetric Laplacian
$L$ has diagonal $d_i$ and off-diagonal $-1$ at neighbouring nodes. Each
undirected connection contributes once to this Laplacian; the pair of
directed receiving channels does not double its weights.

The physical phase is $\theta_i$, the intrinsic frequency is
$\omega_i^u$, the applied correction is $u_i$, and the requested correction
is $q_i$. The target occupancy of receiving edge $j\to i$ is
$\beta^*_{j\to i}$. The fixed coefficient scale in the native PI law is
$h=p/\omega_{\rm base}$, where $p$ is a local polling period and
$\omega_{\rm base}$ is the implementation reference frequency; $H$ is the
physical sampling period imposed only in this reference. The integral
coefficient is $\kappa_I=k_i h$. These two time scales are kept distinct.
All receiving buffers have common capacity $C>0$, and the intrinsic
frequencies $\omega_i^u$ are constant. The all-ones column vector is $\mathbf1$.

The transpose is denoted by $\top$. A symmetric matrix $A\succeq0$ is
positive semidefinite; $A\succ0$ is positive definite. The matrix order
$A\preceq B$ means $B-A\succeq0$. The symbol $\otimes$ denotes the
Kronecker product, and $\|z\|_\infty=\max_i|z_i|$.

\subsubsection{Compatible occupancy targets and marking}
Throughout the proved centered reference, all initial phases are equal.
For zero delay the fluid law in the additive phase-offset convention is
\begin{equation}
 \beta_{j\to i}(t)=\theta_j(t)-\theta_i(t)+\lambda_{j\to i}.
 \label{eq:btd_d_fluid}
\end{equation}
Here $\lambda_{j\to i}$ is the fixed additive phase offset. With equal
initial phases it also equals the initial marking; this presentation does
not identify these conventions for arbitrary unequal initial phases.
The law follows the bittide fluid modelling framework
\cite{lall2024logicalSynchrony,lall2022resistance}. Define
$\beta^*_{j\to i}=\theta_j(0)-\theta_i(0)+\lambda_{j\to i}$, so the
reference starts at its compatible targets. For a common constant ramp
frequency $\omega_{\rm ref}$ put
\begin{equation}
 x_k=\theta(t_k)-\theta(0)-t_k\omega_{\rm ref}\mathbf1.
 \label{eq:btd_ref_state}
\end{equation}
The common ramp cancels from edge differences. Subtracting the target in
\eqref{eq:btd_d_fluid} gives
\begin{equation}
 \beta_{j\to i}(t_k)-\beta^*_{j\to i}=x_{j,k}-x_{i,k},\qquad E_k=-Lx_k.
 \label{eq:btd_buffer_phase}
\end{equation}
The vector $E$ collects the Section~\ref{sec:fpi_controller} frame errors, actual minus target.
Not every set of per-edge targets is compatible with a fixed marking and
phase initialization; that compatibility is an assumption of this model.

For the worked example, equal initial phases and target 1024 imply
$\lambda_{j\to i}=1024$. The actual delayed runs have marking 1044 with
20 initial in-flight frames; simply retaining 1044 while imposing zero
delay and the same centered occupancy would change the initialization.

\subsection{From the original controller to the sampled recurrence}

Define the intrinsic offsets $\delta\omega=\omega^u-\omega_{\rm ref}\mathbf1$
and the physical frequency offsets $w_k=\delta\omega+u_k$. For $k\ge1$,
let $v_k=Hw_{k-1}$. The plant's previous genuine interval gives
\begin{equation}
 x_k-x_{k-1}=Hw_{k-1}=v_k.
 \label{eq:btd_d_plant}
\end{equation}
Indeed, on each zero-order-held interval,
\begin{equation}
 x(t)=x_k+(t-t_k)w_k,\qquad t\in[t_k,t_{k+1}],
 \label{eq:btd_exact_continuous_flow}
\end{equation}
where $x(t)=\theta(t)-\theta(0)-t\omega_{\rm ref}\mathbf1$.
Integration gives $x_{k+1}=x_k+Hw_k$ exactly. The controller is updated
only at the endpoints, but the plant evolves continuously throughout
the interval.

The factor $H$ is physical elapsed time; the controller factor $h$ remains
the native fixed coefficient scale. Section~\ref{sec:fpi_controller} uses
$e_i=-\gamma_i E_i$, so the Fuzzy inputs are
\begin{equation}
 X_k=K_e\Gamma Lx_k,\qquad
 Y_k=K_{\Delta e}\Gamma Lv_k,\qquad
 \Gamma=\operatorname{diag}(\gamma_1,\ldots,\gamma_n).
 \label{eq:btd_d_inputs}
\end{equation}
Here $K_e,K_{\Delta e}$ are chosen common for readability; node-dependent
output scaling can still produce common calibrated $k_p,k_i$.

On $\Omega_r=\{(X,Y):|X|+|Y|\le r\}$, $0<r\le5$, Section~\ref{sec:fpi_controller} gives
$F=\alpha F_0$, where $F_0=(X+Y)/2$ and
$\underline\alpha=(1+r/5)^{-1}\le\alpha\le\overline\alpha=2$.
With $D_{\alpha,k}=\operatorname{diag}(\alpha_{1,k},\ldots,\alpha_{n,k})$,
the calibrated increment becomes
\begin{equation}
 u_k-u_{k-1}=-D_{\alpha,k}(\kappa_I Lx_k+k_pLv_k).
 \label{eq:btd_d_command_increment}
\end{equation}
Multiplication by $H$ and the next plant interval yield
\begin{align}
 v_{k+1}&=v_k-D_{\alpha,k}(bLx_k+aLv_k),\nonumber\\
 x_{k+1}&=x_k+v_{k+1},\qquad
 a=Hk_p,\quad b=H\kappa_I=Hk_i h.
 \label{eq:btd_full_recurrence}
\end{align}
The extra $H$ in $b$ is essential: $v$ is a scaled frequency. The surface
factor $\alpha_{i,k}$ is node-dependent and state-dependent; freezing it
and checking one constant-gain PI characteristic is insufficient.

\subsubsection{The initial constructor call}
The controller's stored previous error is not a previous plant interval.
For the centered initialization $x_0=0$ with zero PI sample sum, zero
Fuzzy accumulator and zero previous percentage error, the first original
controller call gives $u_0=0$. Advancing $[0,H]$ then gives
\begin{equation}
 x_1=H\delta\omega,\qquad v_1=H\delta\omega.
 \label{eq:btd_bootstrap}
\end{equation}
Only from $k=1$ does~\eqref{eq:btd_d_plant} justify the recurrence.
Nonzero initial occupancy error is outside this centered formulation:
one must also introduce target-compatible phase coordinates or retain the
affine error bias, then explicitly evaluate the original constructor and
recheck the initial state and certificate. Evaluating the first call alone
does not restore $E=-Lx$. The centered bootstrap is consistent with, but separate from, the
event-level first-call PI/$F_0$ equivalence proved in Section~\ref{sec:fpi_controller}.

\subsection{Disagreement projection and the symmetric sector matrix}

Since $L\mathbf1=0$, a common phase shift does not change occupancies.
Choose $U$ with orthonormal columns spanning the subspace orthogonal to
$\mathbf1$, so $U^\top U=I_{n-1}$ and $U^\top\mathbf1=0$. Choose its
columns as graph eigenvectors, giving
\begin{equation}
 L=U\Lambda U^\top,\qquad
 \Lambda=\operatorname{diag}(\mu_2,\ldots,\mu_n),\quad 0<\mu_2\le\mu_n.
 \label{eq:btd_d_spectrum}
\end{equation}
$I_{n-1}$ denotes the identity matrix of dimension $n-1$. Define
\begin{equation}
 \xi_k=\Lambda^{1/2}U^\top x_k,\qquad
 \eta_k=\Lambda^{1/2}U^\top v_k,\qquad y_k=(\xi_k^\top,\eta_k^\top)^\top.
 \label{eq:btd_projection}
\end{equation}
The symbol $\xi$ is projected phase and $\eta$ already contains the
frequency scale $H$. The coordinates retain the physical-time scale $H$; in particular,
$\eta$ is a scaled frequency disagreement rather than an integral-controller state.

For $M_k=\Lambda^{1/2}U^\top D_{\alpha,k}U\Lambda^{1/2}$,
substitution in~\eqref{eq:btd_full_recurrence} gives
\begin{align}
 \eta_{k+1}&=-bM_k\xi_k+(I_{n-1}-aM_k)\eta_k,\nonumber\\
 \xi_{k+1}&=(I_{n-1}-bM_k)\xi_k+(I_{n-1}-aM_k)\eta_k.
 \label{eq:btd_d_projected_recurrence}
\end{align}
In particular, the diagonal $D_\alpha$ need not commute with $L$; the
projection produces the symmetric $M_k$ instead of assuming shared modes.
Its quadratic form satisfies
\begin{equation}
 \underline\alpha\Lambda\preceq M_k\preceq\overline\alpha\Lambda,
 \qquad \operatorname{eig}(M_k)\subseteq
 [\nu_-,\nu_+]=[\underline\alpha\mu_2,\overline\alpha\mu_n].
 \label{eq:btd_d_sector_matrix}
\end{equation}
The scalar $\nu$ used below is a possible eigenvalue of $M$, not a node
frequency or the number of communication edges.

\subsection{A common Lyapunov certificate from two endpoints}
\label{subsec:btd_ref_certificate}

Define
For matrices $M\in\mathbb R^{(n-1)\times(n-1)}$ and scalars $\nu$,
define
\begin{equation}
 \mathcal B(M)=\begin{bmatrix}I-bM&I-aM\\-bM&I-aM\end{bmatrix}.
 \label{eq:btd_matrix_dynamics}
\end{equation}
The identity blocks have dimension $n-1$. Equations above give
$y_{k+1}=\mathcal B(M_k)y_k$. A scalar block associated with an eigenvalue
of $M_k$ is
\begin{equation}
 B(\nu)=\begin{bmatrix}1-b\nu&1-a\nu\\-b\nu&1-a\nu\end{bmatrix}.
 \label{eq:btd_scalar_block}
\end{equation}
For a symmetric positive-definite matrix $\bar P\in\mathbb R^{2\times2}$,
put $P=\bar P\otimes I_{n-1}$. The scalar
$V_k=y_k^\top P y_k$ is a quadratic measure of phase and frequency
disagreement, not of absolute phase.

\begin{btdlemma}[Endpoint certificate]
\label{lem:btd_endpoint}
Suppose $0<\rho<1$ and
\begin{align}
 B(\nu_-)^\top\bar P B(\nu_-)&\preceq\rho\bar P,\nonumber\\
 B(\nu_+)^\top\bar P B(\nu_+)&\preceq\rho\bar P.
 \label{eq:btd_d_endpoints}
\end{align}
Then every update whose inputs satisfy the regional sector obeys
\begin{equation}
 V_{k+1}\le\rho V_k.
 \label{eq:btd_contraction}
\end{equation}
The conclusion permits node-dependent and time-varying sector factors.
\end{btdlemma}
\begin{proof}
First consider one possible scalar eigenvalue. For
$\nu=(1-t)\nu_-+t\nu_+$, $0\le t\le1$, affine dependence gives
$B(\nu)=(1-t)B_-+tB_+$, where $B_\pm=B(\nu_\pm)$.
Expanding the two quadratic expressions yields
\begin{align}
 &(1-t)B_-^\top\bar P B_-+tB_+^\top\bar P B_+
       -B(\nu)^\top\bar P B(\nu)\nonumber\\
 &\qquad=t(1-t)(B_+-B_-)^\top\bar P(B_+-B_-)\succeq0.
 \label{eq:btd_d_convexity}
\end{align}
Together with the endpoint assumptions, this proves
$B(\nu)^\top\bar P B(\nu)\preceq\rho\bar P$ throughout the
continuous interval $[\nu_-,\nu_+]$. No eigenvalue grid is required.

Next, each $M_k$ is real symmetric, so it admits an orthogonal
factorization $M_k=Q_k\operatorname{diag}(\nu_{1,k},\ldots,
\nu_{n-1,k})Q_k^\top$. Rotating both $\xi_k$ and $\eta_k$ by $Q_k^\top$
reduces the dynamics to the scalar blocks $B(\nu_{\ell,k})$.
The corresponding block rotation leaves $P=\bar P\otimes I_{n-1}$
unchanged. Thus summing the scalar inequalities gives
$\mathcal B(M_k)^\top P\mathcal B(M_k)\preceq\rho P$, which is
\eqref{eq:btd_contraction}. Because the witness $P$ is the same for
every permitted $M_k$, $Q_k$ may change at every step. No simultaneous
diagonalization of the node-dependent matrices is assumed.
\end{proof}

The remaining task is to show that the actual nonlinear inputs stay in
the region where the sector was obtained. Stability of a family of
hypothetical linear blocks alone would not establish that fact.

\subsection{Closing the invariant input region}

Let $\mathbf e_i$ denote the $i$th unit vector and define the row
$a_i^\top=\mathbf e_i^\top U\Lambda^{1/2}$. Its squared Euclidean norm
is $\|a_i\|^2=L_{ii}=d_i$. The Fuzzy inputs become
\begin{equation}
 X_{i,k}=\gamma_iK_ea_i^\top\xi_k,\qquad
 Y_{i,k}=\gamma_iK_{\Delta e}a_i^\top\eta_k.
 \label{eq:btd_d_linear_outputs}
\end{equation}
For any vector $c$ and $P\succ0$, Cauchy--Schwarz in the quadratic
metric gives $(c^\top y)^2\le(y^\top P y)(c^\top P^{-1}c)$.
Also $|X|+|Y|=\max_{\sigma=\pm1}|X+\sigma Y|$. Applying that output
bound on $V\le R$ yields
\begin{align}
 |X_{i,k}|+|Y_{i,k}|&\le\sqrt{R\Psi_i},\nonumber\\
 \Psi_i&=\gamma_i^2d_i\max_{\sigma=\pm1}
 [K_e,\sigma K_{\Delta e}]\bar P^{-1}
 [K_e,\sigma K_{\Delta e}]^\top.
 \label{eq:btd_input_factor}
\end{align}
The scalar $\Psi_i$ is the largest squared input gain on the energy
ellipsoid. If $V_1\le R$ and $R\Psi_i\le r^2$ for every node, then:
each node's first scaled input pair belongs to $\Omega_r$; the valid sector gives
$V_2\le\rho V_1\le R$; the same output bound keeps the next inputs in
$\Omega_r$; induction repeats. This proves the sector domain and
disagreement ellipsoid invariant together, rather than assuming either.

Equivalently, a valid radius must satisfy
\begin{equation}
 V_1\le R\le\min_i\frac{r^2}{\Psi_i}.
 \label{eq:btd_sector_radius_condition}
\end{equation}
The initialization, normalization through $d_i$ and $C$, the Fuzzy input
scales and the graph-dependent witness therefore enter the same
condition. When $\Psi_i=0$, that node places no restriction on $R$;
otherwise the displayed quotient is finite. The inference surface is
still nonlinear inside the invariant region. The sector is an enclosure
of its instantaneous response, not an assertion that its gain is fixed.

Because $P\succ0$ and $V_k\le\rho^{k-1}V_1$, the projected phase and
frequency decay exponentially. The ellipsoid constrains disagreement;
the common physical phase continues to grow and is not bounded by $V$.
Thus convergence means agreement of rates and restoration of compatible
occupancy targets, not convergence of the absolute oscillator phases to constants.

\subsection{From disagreement to physical buffer and clock bounds}

\subsubsection{Edge occupancy and intersample time}
For edge $e=(j\to i)$, define its signed unit incidence vector
$\mathbf b_e=\mathbf e_j-\mathbf e_i$. Since it is orthogonal to
$\mathbf1$, its output from the projected phase is
\[
 \mathbf b_e^\top x_k=\mathbf b_e^\top U\Lambda^{-1/2}\xi_k.
\]
Define the graph pseudoinverse $L^\dagger=U\Lambda^{-1}U^\top$ and
effective resistance $R_e^{\rm eff}=\mathbf b_e^\top L^\dagger\mathbf b_e$.
The quadratic output inequality gives
\begin{equation}
 |\beta_e(t_k)-\beta_e^*|^2
 \le V_k[\bar P^{-1}]_{11}R_e^{\rm eff}.
 \label{eq:btd_edge_bound}
\end{equation}
Under ZOH the phase, hence each fluid occupancy, is affine on an interval.
For $0\le s\le1$ its deviation is a convex combination of endpoint
deviations, so its absolute value cannot exceed the larger endpoint
absolute value. This covers all physical time, including the initial
interval from the zero state to $x_1$. It would not apply unchanged to
arbitrary delayed integer arrival events.

\begin{btdcorollary}[Buffer safety within the reference model]
\label{cor:btd_buffer_safety}
Suppose the compatible centered initialization and the invariant-domain
conditions above hold with $V_1\le R$. For each receiving edge define
\begin{align}
 m_e&=\min\{\beta_e^*,C-\beta_e^*\},\nonumber\\
 B_e(R)&=\sqrt{R[\bar P^{-1}]_{11}R_e^{\rm eff}}.
 \label{eq:btd_buffer_headroom_condition}
\end{align}
If $0<\beta_e^*<C$ and $B_e(R)<m_e$ on every receiving edge,
then the unconstrained fluid execution satisfies
$0<\beta_e(t)<C$ for all $t\ge0$. In addition,
$\beta_e(t)-\beta_e^*\to0$ as $t\to\infty$.
\end{btdcorollary}
\begin{proof}
At $t_k$, the squared deviation is bounded by
$V_k[\bar P^{-1}]_{11}R_e^{\rm eff}$ and hence by $B_e(R)^2$.
The initial deviation is zero. On $[t_k,t_{k+1}]$, the deviation is a
convex combination of its endpoint deviations, so its absolute value
is at most $B_e(R)$. The strict headroom inequalities give the lower
and upper capacity inequalities. Finally $V_k\le\rho^{k-1}V_1\to0$
and the same interpolation argument imply convergence at every
physical time, not only at samples.
\end{proof}

Here $B_e(R)<m_e$ is a sufficient condition. If it fails, the certificate
is inconclusive; failure of this conservative test does not prove a
buffer overflow or underflow. The condition also concerns the
unconstrained flow. It does not identify a protection-induced stall
with successful control convergence.

\subsubsection{Frequency-difference safety and convergence}
Define the physical frequency spread by
\begin{equation}
 \Delta_\omega(t)=\max_i\omega_i(t)-\min_i\omega_i(t).
 \label{eq:btd_frequency_spread}
\end{equation}
A small frequency spread means close oscillator rates; it does not
require the phases or their integer counts to coincide. For any pair
$i\ne j$, not necessarily adjacent, define
\begin{equation}
 R_{ij}^{\rm eff}=(\mathbf e_i-\mathbf e_j)^\top L^\dagger
                     (\mathbf e_i-\mathbf e_j),
 \qquad R_{\max}^{\rm eff}=\max_{i\ne j}R_{ij}^{\rm eff},
 \label{eq:btd_pair_resistance}
\end{equation}
where $\mathbf e_i$ is the $i$th unit vector. Connectedness makes these
quantities finite. During $[t_k,t_{k+1})$, $\omega_i(t)=
\omega_{\rm ref}+w_{i,k}$, while $v_{k+1}=Hw_k$. Therefore
\begin{align}
 |\omega_i(t)-\omega_j(t)|^2
 &=H^{-2}|(\mathbf e_i-\mathbf e_j)^\top v_{k+1}|^2\nonumber\\
 &\le H^{-2}V_{k+1}[\bar P^{-1}]_{22}R_{ij}^{\rm eff}\nonumber\\
 &\le H^{-2}R\rho^k[\bar P^{-1}]_{22}R_{ij}^{\rm eff}.
 \label{eq:btd_pair_frequency_bound}
\end{align}
The factor $V_{k+1}$, rather than $V_k$, is essential: $v_k$ represents
the frequency held during the previous interval. The case $k=0$ is
covered by the explicitly constructed bootstrap state $y_1$.
At switching times the same bound applies to the newly held frequency;
the previous interval bounds the left-hand limit.

\begin{btdcorollary}[A prescribed frequency-spread threshold]
\label{cor:btd_frequency_safety}
Under the invariant-domain conditions, let $\Delta_{\rm safe}>0$ be a
specified engineering tolerance on pairwise frequency differences. If
\begin{equation}
 D(R):=\frac1H\sqrt{R[\bar P^{-1}]_{22}R_{\max}^{\rm eff}}
       <\Delta_{\rm safe},
 \label{eq:btd_frequency_threshold_condition}
\end{equation}
then for every $t\in[t_k,t_{k+1})$, $k\ge0$,
\begin{equation}
 \Delta_\omega(t)\le D(R)\rho^{k/2}<\Delta_{\rm safe},
 \qquad \Delta_\omega(t)\longrightarrow0.
 \label{eq:btd_frequency_decay}
\end{equation}
\end{btdcorollary}
\begin{proof}
The spread equals $\max_{i,j}|\omega_i-\omega_j|$. Maximizing
\eqref{eq:btd_pair_frequency_bound} and taking square roots proves
the bound. Since $0<\rho<1$, the bound tends to zero.
\end{proof}

The threshold is a user-specified performance or hardware tolerance;
logical synchrony does not by itself prescribe its numerical value.
This proof bounds the transient as well as the limiting disagreement.
It neither proves a smaller Fuzzy frequency peak than PI nor adds a
guarantee for arbitrary delayed or stochastic executions. Frequency
spread also cannot certify positive absolute frequency: a collection of
equal nonpositive rates would have zero spread. The common frequency
must therefore be bounded separately.

\subsubsection{The unprojected common frequency}
Let $\bar w_k=n^{-1}\mathbf1^\top w_k$ be the common-frequency offset.
Put $\alpha_c=(\overline\alpha+\underline\alpha)/2$ and
$s_\alpha=(\overline\alpha-\underline\alpha)/2$. With
$g=(\kappa_I,k_p)^\top$, averaging~\eqref{eq:btd_d_command_increment}
and using $\mathbf1^\top U=0$ gives
\begin{equation}
 \bar w_k-\bar w_{k-1}
 =-\frac1n\mathbf1^\top(D_{\alpha,k}-\alpha_c I_n)
 U\Lambda^{1/2}(\kappa_I\xi_k+k_p\eta_k).
 \label{eq:btd_d_common_increment}
\end{equation}
The second component of $g$ is $k_p$, not $Hk_p$, since $\eta$ already
contains $H$. Bounding each factor yields
\begin{equation}
 |\bar w_k-\bar w_{k-1}|
 \le s_\alpha\sqrt{\frac{\mu_n}{n}g^\top\bar P^{-1}g}\sqrt{V_k}.
 \label{eq:btd_d_common_step_bound}
\end{equation}
The geometric series in $\sqrt{\rho}$ is summable. Hence
\begin{align}
 \sup_k|\bar w_k|
 &\le |\bar w_0|+
 \frac{s_\alpha\sqrt{(\mu_n/n)R g^\top\bar P^{-1}g}}{1-\sqrt\rho}
 \nonumber\\
 &\le |\bar w_0|+\frac{2s_\alpha}{1-\rho}
 \sqrt{\frac{\mu_n}{n}R g^\top\bar P^{-1}g}=W_c.
 \label{eq:btd_clock_bounds}
\end{align}
The last inequality follows from
$(1-\sqrt\rho)^{-1}=(1+\sqrt\rho)/(1-\rho)\le2/(1-\rho)$.
Absolute summability also proves convergence of $\bar w_k$. Combined with
vanishing disagreement, all node frequencies converge to one value
$\omega_\infty=\omega_{\rm ref}+\lim_k\bar w_k$.
For Fuzzy PI that limit need not be the initial intrinsic mean.

\subsubsection{Closing actuator saturation and positive frequency}
The $i$th disagreement-frequency output satisfies
\begin{equation}
 |w_{i,k}-\bar w_k|\le W_d,
 \qquad W_d=\max_i\frac1H
 \sqrt{R[\bar P^{-1}]_{22}(L^\dagger)_{ii}}.
 \label{eq:btd_d_node_frequency}
\end{equation}
Applying the state bound to $y_{k+1}$ covers $w_k$; $y_1$ covers $w_0$.
As $u=w-\delta\omega$, we obtain
\begin{equation}
 |u_{i,k}|\le W_c+W_d+\|\delta\omega\|_\infty,
 \qquad \omega_i(t)\ge\omega_{\rm ref}-W_c-W_d.
 \label{eq:btd_actuator_bounds}
\end{equation}
First derive these bounds for the unconstrained trajectory. If the first
is strictly below $u_{\max}$ and the second is positive, the limiter is
never reached, so the constrained execution coincides with that trajectory.
This continuation argument closes the assumption of no clipping. The
buffer bound must likewise fit within physical capacity; protecting a
buffer by blocking execution is a different outcome.

Together, the following sufficient conditions connect the controller
and network parameters to the required properties:
\begin{align}
 &V_1\le R,\qquad R\Psi_i\le r^2\quad\text{for every }i,\nonumber\\
 &0<\beta_e^*<C\quad\text{for every }e,\nonumber\\
 &R[\bar P^{-1}]_{11}R_e^{\rm eff}<m_e^2
       \quad\text{for every }e,\nonumber\\
 &R[\bar P^{-1}]_{22}R_{\max}^{\rm eff}
       <H^2\Delta_{\rm safe}^2,\nonumber\\
 &W_c+W_d+\|\delta\omega\|_\infty<u_{\max},\nonumber\\
 &\omega_{\rm ref}-W_c-W_d>0,
 \label{eq:btd_joint_sufficient_conditions}
\end{align}
provided the endpoint inequalities of Lemma~\ref{lem:btd_endpoint} hold.
The frequency-spread inequality is required only when a specific threshold
is an objective; convergence alone does not require choosing that threshold.
All quantities are evaluated for one consistent model, initialization
and parameter set. In particular, changing $H$, $h$, $C$, node degree,
controller gains or topology can change more than one inequality.
These are sufficient certification tests, not a characterization of the
largest safe set or necessary stability boundaries.

\subsection{The prescribed triangle certificate}

Take the three-node complete graph, giving
$\mu_2=\mu_3=3$, $R_e^{\rm eff}=2/3$ and
$(L^\dagger)_{ii}=2/9$. The parameters are
\begin{align}
 H=h&=\frac{30000000}{29999},&k_p&=\frac1{500000},&k_i&=\frac1{20000000000},\nonumber\\
 K_e&=\frac1{10},&K_{\Delta e}&=\frac{29999}{7500},&K_u&=-\frac{768}{18749375},\nonumber\\
 C&=2048,&\beta^*_e&=\beta_e(0)=1024,&\lambda_e&=1024.
 \label{eq:btd_example_gains}
\end{align}
Frequencies in this numerical example are normalized, and physical time is measured in reference ticks. Use $\omega^u=(1.0002,0.9998,0.9999)^\top$ and its average as
$\omega_{\rm ref}=29999/30000$, so
$\delta\omega=(7,-5,-2)^\top/30000$ and $\bar w_0=0$.
At the compatible startup~\eqref{eq:btd_bootstrap}, use
\begin{equation}
 \bar P=\frac1{50000}\begin{bmatrix}9&181\\181&49629\end{bmatrix},
 \quad\rho=\frac{1599}{1600},\quad r=1.
 \label{eq:btd_witness}
\end{equation}
Here $\det\bar P=4139/25000000>0$ and the leading diagonal entry is
positive. Since the sum of its four entries is one,
\begin{equation}
 R_0=V_1=H^2\delta\omega^\top L\delta\omega
     =\frac{234000000}{899940001}.
 \label{eq:btd_d_R0}
\end{equation}
This exact fraction is used in all inequalities. The displayed approximation
$0.2600173342$ is slightly below the exact value and is not substituted as
an upper energy radius.

The sector is $[5/6,2]$, hence the spectral endpoints are $5/2$ and 6.
For $D_\nu=\rho\bar P-B(\nu)^\top\bar P B(\nu)$, exact rational
arithmetic verifies positive leading minor and determinant at both ends:
\begin{center}\small
\begin{tabular}{@{}lll@{}}\toprule
$\nu$ & $[D_\nu]_{11}$ (approx.) & $\det D_\nu$ (approx.)\\\midrule
$5/2$ & $8.2193625299\times10^{-7}$ & $7.8876135318\times10^{-11}$\\
$6$ & $2.0776400066\times10^{-6}$ & $7.3380136479\times10^{-9}$\\
\bottomrule\end{tabular}\end{center}
The accompanying machine-readable certificate stores the exact fractions;
the decimals are only a readable display, not rounded tests of positivity.
Using $R_0$ gives the following conservative, outward bounds:
\begin{center}\small
\begin{tabular}{@{}ll@{}}\toprule
Quantity & Certified bound\\\midrule
$\sup_{i,k}(|X_{i,k}|+|Y_{i,k}|)$ & $<0.171455<1$\\
Coarse edge excursion & $<32.237442$ frames\\
$\sup_k|\bar w_k|$ & $<0.003682181$\\
$\sup_{i,k}|w_{i,k}-\bar w_k|$ & $<0.000250634$\\
$\sup_{i,k}|u_{i,k}|$ & $<0.004166148<0.01$\\
$\inf_{i,t}\omega_i(t)$ & $>0.99603385$\\
\bottomrule\end{tabular}\end{center}
The coarse excursion already fits within the 1024-frame lower and upper
headroom. The frequency bounds use the exact rational-square-root
enclosures, not a sum of prematurely rounded display values.

\subsection{The exact PI all-time baseline}

For $J\in\{\mathrm{PI},\mathrm F\}$ define
\begin{equation}
 B_J=\sup_{t\ge0}\max_e|\beta_e^J(t)-1024|.
 \label{eq:btd_excursion}
\end{equation}
The supremum varies time and receiving edges, while the initialization
and intrinsic-frequency vector remain fixed. For the same plant, initialization and gains, native PI has
$D_\alpha=I_n$. Averaging its increments shows that the mean frequency is
preserved. On the triangle, both disagreement modes have eigenvalue 3.
Define $\chi_0=0$, $\chi_1=1$ and
\begin{equation}
 \chi_{k+1}=(2-3a-3b)\chi_k-(1-3a)\chi_{k-1}.
 \label{eq:btd_pi_scalar}
\end{equation}
To derive it, use $v_k=x_k-x_{k-1}$ in the PI state recurrence:
\[
 x_{k+1}=(2I_n-(a+b)L)x_k-(I_n-aL)x_{k-1}.
\]
Since $L\delta\omega=3\delta\omega$ and $x_1=H\delta\omega$,
induction gives $x_k=H\chi_k\delta\omega$. With
$\Delta_0=\max_i\omega_i^u-\min_i\omega_i^u=1/2500$,
\begin{equation}
 \max_e|\beta_e(t_k)-1024|=H\Delta_0|\chi_k|.
 \label{eq:btd_d_pi_deviation}
\end{equation}
This exact rational sequence is enumerated through $N=200$. Its unique
prefix maximum occurs at $k=111$, with
\begin{equation}
 23.40370130812701\le
 \max_{0\le k\le N}H\Delta_0|\chi_k|\le23.40370130812702.
 \label{eq:btd_d_pi_peak}
\end{equation}
It is not yet an all-time assertion until the tail is bounded.

At $N=200$ put $\delta\chi_N=\chi_N-\chi_{N-1}$, the scalar response difference. The exact state energy is
\begin{equation}
 V_N=R_0\{\bar P_{11}\chi_N^2+2\bar P_{12}\chi_N\delta\chi_N+\bar P_{22}(\delta\chi_N)^2\}.
 \label{eq:btd_d_pi_VN}
\end{equation}
The common certificate also covers PI's scalar value $\nu=3$. Thus for
all $k\ge N$, the remaining buffer excursion is at most
$\sqrt{(2/3)[\bar P^{-1}]_{11}V_N}$, with an outward upper bound
$17.95805698665804$ frames. This is strictly below the attained prefix
maximum. ZOH interpolation therefore makes~\eqref{eq:btd_d_pi_peak}
the exact all-physical-time $B_{\rm PI}$ enclosure, reached at time $111H$.

\subsection{The Fuzzy interval prefix and contracting tail}

A direct floating-point trajectory can locate a candidate peak but does
not enclose the exact one. The certificate instead propagates intervals
for all components of $x_k$ and $v_k$ on a grid $S^{-1}=10^{-60}$.
A rational number $q_{\rm rat}$ is enclosed by
\[
 \left[\frac{\lfloor Sq_{\rm rat}\rfloor}{S},\frac{\lceil Sq_{\rm rat}\rceil}{S}\right].
\]
Each multiplication uses the extrema of the four endpoint products,
rounded outward; division uses a strictly positive enclosed denominator
and outward-rounded endpoint quotients. Addition and subtraction use
endpoint arithmetic. These operations use integers and rational comparisons,
not floating-point rounding assumptions.

For the central formula, $m=\min(|X|,|Y|)$ is enclosed by interval
absolute values and endpoint minima. Its interaction term can be written
$J=\min(X,Y)$ for two positive inputs, $J=\max(X,Y)$ for two negative
inputs and zero otherwise. It is coordinatewise nondecreasing, including
across the axes; its lower and upper corner values enclose it on any
input rectangle. The positive denominator $1+2m/5$ then permits a valid
interval evaluation of the exact map. The invariant-domain proof ensures
that this central formula is applicable to the exact trajectory.

Propagating the genuine bootstrap state through state $N=200$ gives a
prefix maximum enclosed by
\begin{equation}
 [17.314056817962,\;17.314056817963]\quad\text{frames}.
 \label{eq:btd_d_fuzzy_prefix}
\end{equation}
The prefix retains all node coordinates, so common drift is not silently
discarded. On the triangle the quadratic energy is equivalently
\begin{equation}
 V_k=\sum_{i<j}\left\{
 \bar P_{11}(x_{i,k}-x_{j,k})^2+
 2\bar P_{12}(x_{i,k}-x_{j,k})(v_{i,k}-v_{j,k})+
 \bar P_{22}(v_{i,k}-v_{j,k})^2\right\}.
 \label{eq:btd_d_pair_energy}
\end{equation}
Evaluating this expression outward gives $\overline V_N\ge V_N$.
For all later samples,
\begin{equation}
 |\beta_e(t_k)-1024|^2
 \le\frac23[\bar P^{-1}]_{11}\rho^{k-N}\overline V_N.
 \label{eq:btd_d_fuzzy_tail}
\end{equation}
For either controller, an equivalent tail bound that makes the
finite-prefix/infinite-tail decomposition explicit is
\begin{equation}
 T_N=\sqrt{\frac23[\bar P^{-1}]_{11}\overline V_N}.
 \label{eq:btd_tail}
\end{equation}
For PI one can take $\overline V_N=V_N$ exactly; for Fuzzy,
$\overline V_N$ is an outward interval upper bound. Once $T_N$ is below
the certified prefix lower bound, no later sample can exceed the prefix
maximum. This closes the infinite time horizon without selecting an
arbitrarily long simulation duration.

An outward upper bound for this whole tail is
$12.29271005372874$ frames, strictly below the prefix lower bound.
Including $x_0=0$ and intersample interpolation gives
$B_{\rm F}<17.315$ on the entire physical time axis. The exact Fuzzy
peak is more tightly enclosed by~\eqref{eq:btd_d_fuzzy_prefix}; the
rounded 17.315 value is used for a simple, conservative comparison.

The enclosures used in the comparison are collected in
Table~\ref{tab:btd_prefix_tail}. The prefix intervals are outward
enclosures of actual peak values; the tail entries are upper bounds.
Their different roles are essential for the strict comparison.
\begin{table}[t]
\caption{Certified excursions in the prescribed fluid reference (frames).}
\label{tab:btd_prefix_tail}
\centering\small
\begin{tabular}{@{}lll@{}}
\toprule
Controller & Prefix peak enclosure & Tail upper bound\\
\midrule
PI & $[23.40370130812701,$ & $17.95805698665804$\\
   & $\phantom{[}23.40370130812702]$ & \\
Fuzzy & $[17.314056817962,$ & $12.29271005372874$\\
      & $\phantom{[}17.314056817963]$ & \\
\bottomrule
\end{tabular}
\end{table}

\subsection{Strict comparison, engineering meaning and scope}
\label{subsec:btd_comparative}

\begin{btdtheorem}[Strict headroom benefit for a prescribed initialization]
\label{thm:btd_headroom}
For the triangle, parameters and compatible initialization above,
initialize the PI sample sum, the Fuzzy accumulator and the stored
previous percentage error to zero. Both reference executions keep
positive physical frequencies, remain strictly inside the buffer
interval $(0,2048)$ and never activate the $0.01$ correction limiter.
Buffer deviations and all pairwise frequency differences converge to
zero. Their all-time excursions satisfy
\begin{equation}
 B_{\rm F}<17.315<23.40370130812701\le B_{\rm PI}.
 \label{eq:btd_strict_comparison}
\end{equation}
For the minimum lower/upper headroom $M_J=1024-B_J$,
\begin{equation}
 M_{\rm F}>1006.685,\qquad
 M_{\rm F}-M_{\rm PI}>6.08870130812701\quad\text{frames}.
 \label{eq:btd_margin_improvement}
\end{equation}
\end{btdtheorem}
\begin{proof}
The endpoint inequalities and the radius test certify the invariant
nonlinear input region. The buffer, correction and absolute-frequency
bounds establish capacity safety, absence of clipping and positive
frequencies. The invariant disagreement energy contracts, so the buffer
and frequency-difference conclusions follow. PI is included by taking
$\alpha_i=1$, which belongs to the same sector; PI also preserves its
initial mean frequency. The certified finite-prefix enclosures and the
tail bounds imply that both all-time excursions are attained in the
prefix. The Fuzzy prefix upper bound is below $17.315$, whereas PI
attains the enclosed peak at sample 111. Intersample affine occupancy
rules out a missed larger fluid peak. Finally, each edge has a reverse
edge with the negative occupancy deviation, so the worst lower/upper
headroom is exactly $1024-B_J$. Subtracting the strict excursion bounds
gives the stated gain in headroom.
\end{proof}

The comparison uses a certified Fuzzy upper bound and an attained PI
peak. Merely comparing two conservative upper bounds would not establish
the strict ordering of the actual excursions.

Both reserves are already large. The comparison establishes a conditional
resource-margin benefit, not a demonstrated storage bottleneck or hardware
cost reduction at this operating point.

The supremum is over time and edges for one prescribed intrinsic-frequency
vector. It is not an expanded safe set for all 400-ppm patterns, a guarantee
against arbitrary latency perturbations or superiority over independently
tuned PI. For a different capacity, $\gamma_i=100/(d_iC)$ changes; keeping
the Fuzzy gains unchanged generally changes the token-domain nonlinear
map. A new capacity certificate requires a fresh derivation or a verified
map-preserving normalization change. Physical occupancy constraints also
do not replace the delayed runtime's in-flight and credit/readiness rules.

This is a computer-assisted mathematical certificate with reproducible
integer/rational checks. It is not proof-assistant verification of the
Julia implementation. The actual asynchronous, delayed, quantized cases
belong in the experimental section. This division preserves the distinction
between a theorem about the declared reference and empirical evidence
about Callisto.


% Updated from completed, frozen experiments on 2026-11-03.
% Historical evaluation text is preserved as comments below the active section.
\section{Evaluation}
\label{sec:evaluation}

The evaluation connects the controller analysis to three distinct questions:
whether the continuous-clock implementation respects the PI calibration;
whether the nonlinear reference trajectories agree with the mathematical
analysis; and whether the proposed controller offers a measurable benefit
against independently tuned PI under shared delivery disturbances. We
separate matched-gain mechanism checks from equal-budget controller
selection, and reference-model certificates from actual Callisto measurements.

\subsection{Experimental design and continuous-clock execution}
\label{subsec:eval_design}

All actual experiments use the continuous-phase Callisto realization of
Section~\ref{sec:fpi_controller}. Oscillator phases evolve continuously,
integer frame crossings determine buffer events, and measurements and
correction applications are scheduled by each node's local phase. There
is no common physical polling grid. Table~\ref{tab:eval_setup} gives the
nominal plant; the startup stress study changes the intrinsic frequencies
and reconstructs the associated negative-time histories explicitly.

\begin{table*}[t]
\caption{Nominal actual-Callisto configuration. Time is in physical reference
ticks except where local ticks are explicitly stated. Controller parameters
selected for independent tests are listed separately in
Table~\ref{tab:eval_selected_gains}.}
\label{tab:eval_setup}
\centering
\begin{tabular}{@{}ll@{}}\toprule
Parameter & Value\\\midrule
Topology & Three-node bidirectional triangle; six receiving buffers\\
Intrinsic frequencies $\omega^u$ & $(1.0002,0.9998,0.9999)^\top$; initial spread 400 ppm\\
Initial phase; base frequency & $0.1$ at every node; $\omega_{\rm base}=29999/30000$\\
Capacity; initial occupancy; target & $C=2048$; $\beta_e(0)=\beta_e^*=1024$ frames\\
Initial in-flight frames; marking & 20 per channel; 1044 in the recorded convention\\
Nominal path delay & 20\\
Polling; control-execution delay & $p=1000$; $d_{\rm ctrl}=10$ local ticks\\
Execution constraints & $|u_i|\le0.01$; guards enabled; reserve zero; no anti-windup\\
Control law & Native integrate-first PI or the complete incremental 25-rule Fuzzy law\\
Other mechanisms & Observer and supervisory mode switching disabled\\
Disturbance onset & $t_{\rm on}=2\times10^6$\\\bottomrule
\end{tabular}
\end{table*}

\paragraph{Evidence groups and data separation.}
The implementation check, M0, contains six matched pairs, or 12 actual
runs. An earlier fixed-preset burst campaign, M2, supplies a separate
mechanism diagnostic. The earlier fixed-preset random campaign, M3, is
development evidence and supplies no seeds to the new independent tests.
Its fixed-preset comparison did not support lower Fuzzy RMS against
moderate PI; the new study asks the different question of independently
selected controllers under an equal evaluation budget.
The new training study charges 192 slots, including four verified cached
M2 evaluations and 188 new actual runs. Selection is frozen before 186
new held-out runs. A separately registered slower-disturbance study adds
107 actual runs, including its engineering pilots, and the startup
study adds six. Its 32 synchronous fluid-reference trajectories are
reference calculations, not additional actual runs. Old data, infeasible
candidates and repair records are retained in the evidence archive.

\paragraph{Shared delivery inputs.}
All controllers receive identical externally generated desired-latency
profiles within each comparison, verified by content hashes. Delay is
20 before onset, including the negative-time prehistory. Ordered arrivals
use the causal running-maximum envelope of
Section~\ref{sec:fpi_controller}; a delay decrease may release an ordered
batch. The model does not impose a finite serialization rate.
For bounded random latency, the desired delay is $20+z_m$, where
\begin{align}
 z_0&=5\varepsilon_0,\nonumber\\
 z_m&=r_z z_{m-1}+(1-r_z)5\varepsilon_m,
 \quad r_z=e^{-1},\quad\varepsilon_m\sim\mathcal U[-1,1].
 \label{eq:eval_random_profile}
\end{align}
Thus desired delay remains in $[15,25]$. The index $m$ counts physical
profile updates on a grid of spacing $T_c$, independently of controller
polling. Here $T_c$ denotes a profile-update interval, not an exactly
measured correlation time. The original two regimes use $T_c=1000$
and $10000$; the later confirmation study uses $100000$.

\subsection{Metrics, admissibility and comparison windows}
\label{subsec:eval_metrics}

For $n$ nodes, let $\bar\omega(t)=n^{-1}\sum_i\omega_i(t)$ and define
the instantaneous across-node root-mean-square (RMS) disagreement as
\begin{align}
 \sigma_\omega(t)&=\sqrt{n^{-1}\sum_i[\omega_i(t)-\bar\omega(t)]^2},\nonumber\\
 \sigma_{\rm ppm}(t)&=10^6\sigma_\omega(t)/\omega_{\rm unit},
 \qquad\omega_{\rm unit}=1.
 \label{eq:eval_disagreement_amplitude}
\end{align}
Parts per million (ppm) use the fixed normalized frequency unit, not
$\omega_{\rm base}$. For a physical-time window $W=[a,b)$ of duration
$T_W=b-a$, report
\begin{align}
 R_W&=\sqrt{T_W^{-1}\int_W\sigma_{\rm ppm}(t)^2\,dt},\nonumber\\
 I_W&=\int_W\sigma_{\rm ppm}(t)\,dt,
 \qquad P_W=\sup_{t\in W}\sigma_{\rm ppm}(t).
 \label{eq:eval_frequency_metrics}
\end{align}
The integral absolute error (IAE) $I_W$ integrates RMS amplitude, rather
than the mean absolute node deviation. Frequency moments use complete
held-frequency intervals weighted by their physical durations.
The actual pairwise spread $\Delta_\omega=\max_i\omega_i-\min_i\omega_i$
is recorded separately. Although
$\Delta_\omega\le\sqrt{2n}\,\sigma_\omega$ pointwise, window RMS alone
does not establish a maximum transient-spread bound.

Buffer excursion, minimum two-sided headroom and mean command total
variation (TV) are
\begin{align}
 B_W^{\rm obs}&=\max_e\sup_{t\in[a,b]}|\beta_e(t)-\beta_e^*|,\nonumber\\
 H_W^{\rm obs}&=\min_e\inf_{t\in[a,b]}
                         \min\{\beta_e(t),C-\beta_e(t)\},
 \label{eq:eval_buffer_metrics}\\
 \mathrm{TV}_W&=n^{-1}\sum_i
       \sum_{k:\,a\le t^a_{i,k}\le b}|u_{i,k}-u_{i,k-1}|.
 \label{eq:eval_tv}
\end{align}
Buffer statistics reconstruct all integer crossings and ordered arrivals,
grouping simultaneous events and including the endpoint state. Sparse
plot samples do not determine extrema. TV includes both endpoints under
the frozen convention, so adjacent-window activities cannot be added
without treating the shared endpoint. TV measures command activity,
not hardware energy. Requested and applied corrections, saturation,
gate events, positive frequency and data-phase progress are checked
separately.

Before each post-centering disturbance, three consecutive $10^5$-tick
windows ending at onset must satisfy disagreement RMS at most 2 ppm,
each buffer's absolute mean bias at most two frames and standard
deviation at most one frame, common-frequency temporal standard
deviation at most 1 ppm and three-window mean range at most 1 ppm,
with positive progress and no intervention. These are finite
engineering centering criteria, not exact equilibrium. A training
candidate failing the same criteria in any required case remains
infeasible and consumes its assigned budget.

The original-regime tests use $W=[2\times10^6,4\times10^6)$ and tail
$[3.6\times10^6,4\times10^6)$. Each controller's own clean trajectory
is its same-parameter, nominal-delay counterfactual over the same
physical times. Common-frequency shift is reported relative to this
counterfactual, separately from relative synchronization. Strict
recovery requires the centering criteria for three consecutive windows
and a window-mean common frequency within 2 ppm of own clean. A return
not observed by the specified endpoint is right-censored, not assigned
a recovery time of zero or interpreted as permanent nonrecovery.

\subsection{Implementation and nonlinear-reference validation}
\label{subsec:eval_implementation_reference}

M0 checks Proposition~\ref{prop:fpi_equivalence} by replacing the Fuzzy
surface with its calibrated linear counterpart $F_0$, without changing
the continuous plant. Six pairs cover polling periods 100 and 1000,
initial frequency spreads 35 and 400 ppm, and capacity/target settings
$(2048,1024)$, $(512,128)$ and $(512,384)$ over $2\times10^5$ ticks.
Across 19,800 paired updates, the maximum correction difference is
$5.29\times10^{-18}$; physical/data phases, measurement vectors and
saved buffer counts agree, without saturation or guard intervention.
This checks initialization, integration order and gain mapping. It does
not replace the complete nonlinear law in the performance experiments.

The original 25-rule law is also evaluated directly on the synchronous,
zero-delay reference of Section~\ref{sec:btd_reference}, with 2000
post-bootstrap updates. At the prescribed moderate-gain initialization,
the fluid buffer peak is approximately 17.314056818 frames, maximum
$|X|+|Y|$ is 0.090130790, maximum absolute correction is
$3.21398377\times10^{-4}$, and no contraction-envelope violation is
observed. Physical frequencies stay in $[0.9998,1.0002]$ including
startup. These finite numerical checks agree with the analysis; the
all-time comparison, Fuzzy below 17.315 versus PI approximately
23.403701 frames, follows from the interval-prefix and contracting-tail
proof in Theorem~\ref{thm:btd_headroom}.

\paragraph{Full-rule and memory checks.}
All 299 new held-out, slower-disturbance and actual startup raw records
pass same-input reconstruction of requested corrections, with maximum
residual approximately $5.42\times10^{-20}$. This uses the complete
rule map and stored accumulator, not a sector surrogate. The earlier
M2 diagnostic additionally stays within $\Omega_5$, with observed
$\max(|X|+|Y|)=2.001888$, supporting application of
Proposition~\ref{prop:fpi_central} to those measured inputs. This
diagnostic is retained separately from the independently selected
controllers: Fuzzy$_B$ reaches approximately 39.159 in the large
held-out burst. The full 25-rule law remains defined there, whereas
the central-region identity and the $\Omega_1$ reference witness do
not cover that whole trace. Same-input memory accounting is an
algebraic diagnostic, not an additional closed-loop $F_0$ experiment.

\subsection{Reference safety strata and actual startup transients}
\label{subsec:eval_safety}

To vary initial mismatch without changing occupancy normalization,
capacity and targets remain fixed while
\begin{equation}
 \omega^u(s)=\frac{29999}{30000}\mathbf1
                 +s\frac{(7,-5,-2)^\top}{30000},
 \quad s\in\{0.5,1,2,4,8,16,32,64\}.
 \label{eq:eval_stress_direction}
\end{equation}
Here $s$ is an initial-mismatch scale, with intrinsic spread $400s$ ppm.
The reference uses the matched moderate coordinates
$(k_p,k_i)=(2\times10^{-6},5\times10^{-11})$, synchronous zero-delay
fluid execution and $H=h=30000000/29999$. Its 32 trajectories comprise
two controllers, eight scales and both unprotected virtual and
$0.01$-clipped variants. Each includes the bootstrap and 2000 subsequent
zero-order-hold intervals. The initial certificate energy scales as
$V_1(s)=s^2V_1(1)$; input-domain, buffer, positive-frequency and correction
conditions are checked jointly.

The largest grid point satisfying the common sufficient conditions is
$s=2$. At $s=4$, both reference trajectories remain normal over the
record although the certificate no longer guarantees them. At $s=32$,
unprotected reference requests first require actuator clipping; at
$s=64$, both virtual reference buffer trajectories violate capacity.
Figure~\ref{fig:eval_reference_grid} distinguishes measured virtual
peaks from sufficient bounds. Failure of a sufficient condition is
inconclusive, rather than a prediction of instability.

\begin{figure*}[t]
\centering
\includegraphics[width=0.97\textwidth]{images/evaluation_latest_20261103/reference_grid.png}
\caption{Synchronous zero-delay fluid-reference sweep at matched moderate
coordinates. Virtual trajectory peaks and analytical sufficient bounds
are distinct quantities. The dashed levels show fluid buffer headroom
and the correction limit; clipping and unprotected virtual trajectories
are archived separately. These are reference calculations, not actual
Callisto safety certificates.}
\Description{Four panels compare reference buffer peaks, requested corrections,
frequency spread and correction bounds over eight initial-mismatch scales.}
\label{fig:eval_reference_grid}
\end{figure*}

Before running actual cases, the selection rule freezes three common
scales: 2, 4 and 32, representing certified, observed-safe uncertified,
and correction-constrained reference strata. Actual Callisto restores
path delay 20 and local polling/control delay 1000/10. Its locked
startup horizon is $2001H\simeq2.001067\times10^6$ physical ticks,
rather than the $4\times10^6$ used for the disturbance tests. This
shorter horizon is an execution-plan deviation from the earlier
completion instruction and limits the empirical safety observation
accordingly. The six cases are startup stress tests, not
post-centering disturbance tests.

\begin{table}[t]
\caption{Actual startup stress at matched moderate coordinates. Buffer
excursions include startup and all reconstructed events. Sat. counts
saturated control updates; every case has zero gate events and replay
mismatches.}
\label{tab:eval_startup_safety}
\centering
\begin{tabular}{@{}rrrrr@{}}\toprule
 & & \multicolumn{2}{c}{Buffer excursion / frames} & Sat.\\
$s$ & Spread / ppm & PI & Fuzzy & PI / Fuzzy\\\midrule
2 & 800 & 47 & 37 & 0 / 0\\
4 & 1600 & 94 & 72 & 0 / 0\\
32 & 12800 & 750 & 664 & 176 / 85\\\bottomrule
\end{tabular}
\end{table}

Fuzzy has a smaller maximum actual buffer excursion at all three
selected scales (Table~\ref{tab:eval_startup_safety}). At $s=4$, the
reduction is from 94 to 72 frames, approximately 23.4\%, with no
saturation or gate intervention. Figure~\ref{fig:eval_startup_actual}
shows the associated frequency and occupancy transients. All six
records retain positive frequency, data progress and capacity
compliance. At $s=32$, Fuzzy also requests fewer saturated updates,
but 12800 ppm is a stress setting, not a typical hardware tolerance.
Its histories are initialized consistently: the faster source has
21 in-flight frames and marking 1045, whereas the other sources retain
20 and 1044. Both controllers receive identical histories at each scale.

\begin{figure*}[t]
\centering
\includegraphics[width=0.97\textwidth]{images/evaluation_latest_20261103/actual_s4.png}
\caption{Actual continuous-phase Callisto startup at $s=4$ (1600 ppm),
with matched moderate coordinates. Upper panels show positive absolute
normalized frequencies using a $+1$ axis offset. Lower panels show the
minimum--maximum occupancy envelope across all receiving buffers, from
the complete crossing stream. Fuzzy reduces the maximum target
deviation from 94 to 72 frames over the stated startup horizon.}
\Description{PI and Fuzzy startup trajectories at the same frequency mismatch,
with absolute node frequencies and complete-event buffer-occupancy envelopes.}
\label{fig:eval_startup_actual}
\end{figure*}

These observations are consistent with a useful transient-headroom
effect under matched gains, but do not transfer the reference theorem
to delayed, asynchronous, quantized execution. The recorded pairwise
spread remains a separate quantity. No hardware-based numerical
$\Delta_{\rm safe}$ is assumed; applying the reference spread corollary
would require an independently justified threshold exceeding $D(R)$.

\subsection{Equal-budget selection and independent random-latency tests}
\label{subsec:eval_random}

\paragraph{Fixed search budget and selection.}
Each family uses the same 24-point grid,
\begin{align}
 k_p&\in\{1.5,2,2.5,3,4,6\}\times10^{-6},\nonumber\\
 k_i&\in\{1,2,5,10\}\times10^{-11}.
 \label{eq:eval_gain_grid}
\end{align}
Fuzzy fixes $K_e=0.1$ and obtains $K_{\Delta e},K_u$ from the
calibration in Section~\ref{sec:fpi_controller}; it receives no extra
shape-search dimension. Four training cases per candidate comprise
clean operation, a single-edge 20-tick delay increment lasting
$10^5$ ticks, and two random regimes using training seed 31001.
Both families therefore receive 96 evaluation slots. Four hash-verified
cached evaluations remain charged within the 192-slot total.
Nineteen PI and 22 Fuzzy candidates satisfy the same finite-run
admissibility requirements; all seven infeasible candidates are
retained without replacement.

Five training objectives are minimized separately: worst random-case
RMS, burst IAE, worst disturbance-window buffer excursion, peak absolute
common-frequency excursion from own clean, and worst mean command TV.
Each family's feasible Pareto set supplies a random-priority
representative, labeled $R$, and a burst-priority representative,
labeled $B$. The frozen lexicographic rule prioritizes random RMS for
$R$ and burst IAE for $B$, with the remaining objectives and candidate
identifier as declared tie breakers. One PI serves both priorities;
Fuzzy selects distinct representatives (Table~\ref{tab:eval_selected_gains}).
Deduplication gives three candidates, without backfilling.

\begin{table*}[t]
\caption{Frozen representatives. The $k_p,k_i$ columns are native PI gains
for PI and calibrated PI coordinates for Fuzzy. $R$ and $B$ denote
random- and burst-priority selection; the candidates are not retuned on tests.}
\label{tab:eval_selected_gains}
\centering
\begin{tabular}{@{}llrrrrr@{}}\toprule
Representative & Candidate & $k_p$ & $k_i$ & $K_e$ & $K_{\Delta e}$ & $K_u$\\\midrule
PI$_R$ = PI$_B$ & PI\_g05 & $2\times10^{-6}$ & $10^{-11}$ & --- & --- & ---\\
Fuzzy$_R$ & Fuzzy\_g02 & $1.5\times10^{-6}$ & $2\times10^{-11}$ & 0.1 & 7.499750 & $-1.638455\times10^{-5}$\\
Fuzzy$_B$ & Fuzzy\_g17 & $4\times10^{-6}$ & $10^{-11}$ & 0.1 & 39.998667 & $-8.192273\times10^{-6}$\\\bottomrule
\end{tabular}
\end{table*}

\paragraph{Untouched paired tests.}
Only after freezing selection are test profiles generated from seeds
41001--41030. All three candidates execute both original random
regimes and both held-out burst combinations, giving
$30\times2\times3+2\times3=186$ new actual runs. All 186 complete
and satisfy the declared normal-admissibility checks, with no
engineering QA errors or complete-event replay mismatches. No seed
pair is excluded. Archived M3 seeds are not reused as untouched tests.

The primary random comparison is Fuzzy$_R$ versus PI$_R$. Let $\ell$
index independent seeds. Report the paired mean difference and mean
of per-seed ratios,
\begin{align}
 \widehat d&=\frac1{30}\sum_{\ell=1}^{30}
       (R_{W,\ell}^{\rm Fuzzy_R}-R_{W,\ell}^{\rm PI_R}),\nonumber\\
 \widehat r&=\frac1{30}\sum_{\ell=1}^{30}
       R_{W,\ell}^{\rm Fuzzy_R}/R_{W,\ell}^{\rm PI_R}.
 \label{eq:eval_paired_ratio}
\end{align}
The ratio is not the ratio of group means. Confidence intervals (CIs)
use 20000 paired-seed percentile-bootstrap resamples, with fixed
resampling seeds 58001/58002. Pointwise 95\% intervals accompany
each endpoint; a joint statement about the two original regimes also
uses nominal Bonferroni-adjusted 97.5\% intervals per regime.
Time points, edges and bootstrap draws are not independent replications.

\begin{table*}[t]
\caption{Untouched original-regime tests: 30/30 admissible pairs per row.
Controller RMS entries are seed means in ppm. Difference is Fuzzy$_R$
minus PI$_R$; ratios are means of paired ratios. Brackets give pointwise
95\% paired bootstrap intervals.}
\label{tab:eval_random}
\centering
\begin{tabular}{@{}lrrll@{}}\toprule
Profile grid & PI$_R$ RMS & Fuzzy$_R$ RMS & Difference / ppm [CI] & Paired ratio [CI]\\\midrule
$T_c=1000$ & 4.61355 & 3.83949 & $-0.77406$ [$-0.80767$,$-0.73992$] & 0.83217 [0.82484,0.83964]\\
$T_c=10000$ & 4.65204 & 4.28718 & $-0.36486$ [$-0.43523$,$-0.29464$] & 0.92150 [0.90630,0.93672]\\\bottomrule
\end{tabular}
\end{table*}

Table~\ref{tab:eval_random} and Figure~\ref{fig:eval_random_primary}
show reductions in disagreement RMS at both tested profile grids.
The mean paired relative improvements are 16.78\% and 7.85\%.
The nominal family-adjusted difference intervals are
$[-0.81274,-0.73483]$ and $[-0.44451,-0.28407]$ ppm, also below zero.
This supports an advantage of the selected Fuzzy scheme under the
declared finite-grid budget, plant and two input regimes. Because
the selected coordinates differ between families, it does not
isolate a same-gain inference-only effect or establish globally
optimal PI/Fuzzy performance.

\begin{figure*}[t]
\centering
\includegraphics[width=0.97\textwidth]{images/evaluation_latest_20261103/paired_random_effects.png}
\caption{Independent paired Fuzzy$_R$/PI$_R$ RMS effects after equal-budget
selection. Each dot is one of 30 untouched seed pairs. The final black
marker is the mean paired ratio with a pointwise 95\% CI; the dashed
line denotes equal RMS. Individual seed outcomes and their aggregate
are both retained.}
\Description{Paired test-seed RMS ratios for profile intervals 1000 and
10000 ticks, with mean effects and bootstrap confidence intervals.}
\label{fig:eval_random_primary}
\end{figure*}

\begin{figure*}[t]
\centering
\includegraphics[width=0.97\textwidth]{images/evaluation_latest_20261103/raw_Tc10000.png}
\caption{Saved raw traces for the first held-out seed (41001) in the
$T_c=10000$ random experiment. Frequency uses held corrections and shared physical
time; buffer occupancy overview samples are for display. Reported
moments and buffer extrema instead use complete intervals and events.
The paired-seed statistics, rather than this one trace, support the
mean performance comparison.}
\Description{Shared-time node-frequency and buffer-occupancy
histories for the first held-out bounded-random delivery input.}
\label{fig:eval_random_raw}
\end{figure*}

\paragraph{Activity and common-frequency cost.}
Mean per-node command TV decreases from 0.0102361 to 0.00782282 at
$T_c=1000$, and from 0.00112361 to 0.000906694 at $T_c=10000$
(PI$_R$ to Fuzzy$_R$). All selected candidates have at most six frames
of disturbance-window target deviation. However, Fuzzy$_R$ has larger
magnitude mean tail common-frequency shifts from own clean:
$-61.420$ versus $-15.692$ ppm, and $-31.933$ versus $-15.384$ ppm.
Thus reduced relative disagreement and command activity do not
imply improved absolute-frequency regulation. No absolute-frequency
reference loop is present.

\subsection{Held-out bursts and slower delivery changes}
\label{subsec:eval_burst}

\paragraph{Burst-priority validation.}
The two predetermined held-out combinations are a 10-tick desired
delay increment lasting $10^4$ ticks on channel $2\to3$, and a
40-tick increment lasting $10^5$ ticks on $3\to1$, both beginning
at $2\times10^6$. They are distinct from the training combination.
For the larger burst, Fuzzy$_B$ lowers post-onset IAE from
$8.19241\times10^6$ to $3.62317\times10^6$ ppm$\cdot$ticks,
approximately 55.8\%, while matching the 40-frame buffer excursion.
Table~\ref{tab:eval_heldout_bursts} retains both deterministic cases:
the improvement depends on the burst combination rather than holding
for every tested burst. These two cases do not supply seed-level CIs.

\begin{table}[t]
\caption{Both held-out burst-priority cases over
$[2\times10^6,4\times10^6)$. IAE is in $10^6$ ppm$\cdot$ticks;
buffer excursion is in frames. $A$ and $D$ denote desired delay
increment and its duration, respectively.}
\label{tab:eval_heldout_bursts}
\centering
\begin{tabular}{@{}lrrrr@{}}\toprule
 & \multicolumn{2}{c}{IAE} & \multicolumn{2}{c}{Buffer excursion}\\
Case & PI$_B$ & Fuzzy$_B$ & PI$_B$ & Fuzzy$_B$\\\midrule
$A=10,D=10^4$ & 0.265814 & 0.352614 & 11 & 11\\
$A=40,D=10^5$ & 8.192410 & 3.623166 & 40 & 40\\\bottomrule
\end{tabular}
\end{table}

\paragraph{Ramp--hold--ramp inputs.}
A separate registered study freezes the selected gains and applies
a 20-tick delay increment on channel $1\to2$ using rise/fall times
$T_r\in\{10^4,10^5\}$ and hold times
$T_h\in\{10^5,10^6\}$. The four combinations all run to
$5.2\times10^6$ ticks. The end of the falling ramp is
$t_{\rm off}=2\times10^6+2T_r+T_h$; each case has its own
$[t_{\rm off},t_{\rm off}+2\times10^6)$ recovery window and
same-horizon clean counterfactual. A larger constant delay plateau
does not mean delivery stops throughout that plateau. The ramp
instead gives feedback time to act while the delivery imbalance forms.

For $T_r=10^5$, maximum buffer deviations are 19 frames for PI,
18 for Fuzzy$_R$ and 17 for Fuzzy$_B$. The slower formation of the
disturbance therefore provides another observed buffer-margin
benefit. Frequency RMS favors PI across the four ramp combinations,
and strict own-clean recovery is censored in all 12 ramp records.
Different disturbance durations are not pooled into a raw-IAE ranking;
these data support a resource/frequency tradeoff, not a faster-settling claim.

\paragraph{Separately registered slow-random confirmation.}
The slower grid $T_c=100000$ uses new seeds 51001--51030, unchanged
selected gains, the same bounded recursion, and
$W=[2\times10^6,5\times10^6)$ with tail
$[4.6\times10^6,5\times10^6)$. Each trace has 30 complete
profile blocks, which are not independent statistical replicates.
The study includes 105 matrix slots: three long-clean cases,
12 ramps and 90 random cases. Four matrix slots reuse qualified
pilot runs; 101 new main runs plus six pilots give 107 unique actual
runs. All matrix slots pass profile, nominal-prefix and event-replay
checks, with no excluded confirmation seed pairs.

This confirmation bounds the applicability of the original-regime
advantage: the mean paired Fuzzy$_R$/PI$_R$ RMS ratio is 1.33312
with pointwise 95\% CI [1.30337,1.36219], using 20000 paired
bootstrap resamples and fixed seed 59001. The mean paired difference
is $+1.53215$ ppm [1.39578,1.66175]. Slower variation alone therefore
does not extend the selected controller's RMS advantage. This confirmation
delimits the two original-regime benefits without changing their
registered primary endpoints or the frozen controller selection.

\subsection{Connection between analysis and measured benefits}
\label{subsec:eval_support}

The evidence establishes two complementary benefits. First, matched-gain
nonlinear analysis gives a strict all-time headroom improvement for the
prescribed fluid reference; actual matched-gain startup tests also show
smaller measured excursions at three fixed initial mismatches. Second,
equal-budget, independent selection yields lower disagreement RMS and
mean command activity in two held-out random regimes, while the
burst-priority representative reduces IAE in the larger held-out burst.
These are distinct statements with different gains, assumptions and
observation windows.

Their scope is explicit: common-frequency shift remains a cost, the
small burst and slower delivery tests delimit performance, and finite
actual records do not certify all-time safety or convergence. The
experiments do not establish a general safe-set inclusion, a globally
optimal comparator, payload-level functional logical synchrony, or
hardware energy savings. Positive phase progress and buffer compliance
show normal model execution; they are not a functional refinement proof.
Within these limits, the combined analysis and tests provide both a
mechanistic reason and independent empirical evidence for considering
the proposed Fuzzy PI law when transient buffer use and
disturbance-dependent synchronization performance matter.


% =====================================================================
% HISTORICAL EVALUATION BEFORE LATEST RESULTS -- ALL LINES COMMENTED
% Preserved for review; old fixed-preset tables/plots are not compiled.
% This archive does not replace the current confirmation-scope summary.
% =====================================================================
% % Evaluation assembled from completed evidence only, 2026-11-03.
% % No new Callisto runs, tuning/selection or PDF compilation is required.
% \section{Evaluation}
% \label{sec:evaluation}
% 
% The evaluation addresses three questions. First, does the implementation
% respect the PI calibration and initialization derived in
% Section~\ref{sec:fpi_controller}? Second, are the reference-model
% calculations consistent with the implemented nonlinear rule law? Third,
% how does the full Fuzzy controller compare with fixed PI baselines under
% post-centering delivery disturbances in continuous-phase Callisto?
% These questions require separate evidence: implementation equivalence,
% reference-model consistency, and finite-horizon performance do not
% establish the same claim.
% 
% \subsection{Evidence groups and comparison design}
% \label{subsec:eval_design}
% 
% We use three completed actual-Callisto campaigns, denoted M0, M2 and M3,
% and separate fluid-reference calculations. M0 comprises six PI/linear-
% counterpart pairs, or 12 actual runs. M2 comprises four actual runs:
% PI and full Fuzzy, each with a clean trace and a single-edge latency
% burst. M3 comprises 180 actual validation runs: 30 independent paired
% seeds, two random-latency regimes and three fixed controllers. Six M3
% pilot runs are separate from the 180-run validation and are not treated
% as independent test replications. Reference calculations are not added
% to the actual-Callisto run count.
% 
% The primary performance comparison is moderate PI against the full
% Fuzzy law calibrated to its PI gain coordinates. M3 additionally includes
% a prespecified stronger PI candidate. All candidates are frozen within
% these campaigns. The linear surface $F_0$ is used only for the M0
% implementation check and explicitly labeled same-input diagnostics;
% M2 and M3 execute the full 25-rule nonlinear law. Axis matching controls
% scales and update conventions, but is not an assertion that either
% candidate is the best independently tuned controller.
% 
% \subsection{Continuous Callisto setup and delivery disturbances}
% 
% M2 and M3 use the continuous oscillator phases, local-clock scheduling
% and integer receiving-buffer counts described in
% Section~\ref{sec:fpi_controller}. Table~\ref{tab:eval_setup} gives their
% common configuration. Reported token counts represent counted frames
% in this model. There is no shared physical controller grid in
% these actual runs. The fixed native integration scale is
% $h=p/\omega_{\rm base}$, with the same convention in the calibrated
% Fuzzy controller. The initial intrinsic-frequency spread is 400 ppm of
% the normalized frequency unit.
% 
% \begin{table*}[t]
% \caption{Common configuration of the aligned burst and random-latency experiments.
% Time is measured in physical reference ticks unless specified as local ticks.}
% \label{tab:eval_setup}
% \centering
% \begin{tabular}{@{}ll@{}}\toprule
% Parameter & Value\\\midrule
% Network & Three-node bidirectional triangle; six receiving buffers\\
% Intrinsic frequencies $\omega^u$ & $(1.0002,0.9998,0.9999)^\top$\\
% Initial phases; implementation base frequency & $0.1$ at every node; $\omega_{\rm base}=29999/30000$\\
% Capacity; initial occupancy; target & $C=2048$; $\beta_e(0)=\beta_e^*=1024$ frames\\
% Initial in-flight frames; delayed-run marking & 20 per channel; 1044 in the recorded convention\\
% Nominal physical path delay & 20\\
% Polling; control-execution delay & $p=1000$; $d_{\rm ctrl}=10$ local ticks\\
% Moderate PI $(k_p,k_i)$ & $(2\times10^{-6},5\times10^{-11})$\\
% Stronger fixed PI, M3 only & $(4\times10^{-6},5\times10^{-11})$\\
% Full Fuzzy $(K_e,K_{\Delta e},K_u)$ & $(0.1,3.9998666667,-4.0961365379\times10^{-5})$\\
% Execution policy & $u_{\max}=0.01$; shared guards enabled; reserve zero; no anti-windup\\
% Optional mechanisms & Observer and supervisory mode switching disabled\\
% Horizon; disturbance onset & $4\times10^6$; $t_{\rm on}=2\times10^6$\\\bottomrule
% \end{tabular}
% \end{table*}
% 
% \paragraph{Burst after nominal centering.}
% In M2, desired send-time latency on one directed channel, identified as
% $1\to2$ in the runtime configuration, changes from 20 to 40 during
% $[2\times10^6,2.1\times10^6)$ and returns to 20 afterwards. All other
% channels remain at 20. Each controller has its own clean counterfactual
% with latency 20 throughout. A controller's burst and clean histories
% match before onset, so their subsequent difference isolates the effect
% of the delivery disturbance for that controller. The nominal states
% of PI and Fuzzy are not required to be identical to each other.
% 
% \paragraph{Bounded random latency.}
% M3 uses fixed physical profile-update grids $T_c=1000$ and $10000$,
% independently of the unchanged controller polling period. For each
% generated link profile the bounded recursion is
% \begin{align}
%  z_0&=5\varepsilon_0,\nonumber\\
%  z_m&=r_z z_{m-1}+(1-r_z)5\varepsilon_m,
%  \qquad r_z=e^{-1},\quad\varepsilon_m\sim\mathcal U[-1,1],
%  \label{eq:eval_random_profile}
% \end{align}
% with desired delay $20+z_m\in[15,25]$ after onset. Here $m$ indexes
% profile updates, not controller measurements. The designation $T_c$
% refers to the profile grid, not an exactly measured correlation time.
% Before onset, including negative-time prehistory, delay is 20. Seeds
% 20001--20030 are independent validation replications; pilot seed 1001
% is excluded. For each seed and regime, all three controllers receive
% identical externally generated profile content.
% 
% Both disturbance families use the declared causal running-maximum
% arrival envelope. Frame order is retained and a delay decrease can
% release an ordered batch. No finite serialization-rate model, frame
% dropout or application-payload experiment is included. These are
% specified transport stress models rather than claims about arbitrary
% network traffic.
% 
% \subsection{Metrics, windows and finite-run admissibility}
% \label{subsec:eval_metrics}
% 
% \paragraph{Disagreement and absolute frequency.}
% For the instantaneous common frequency $\bar\omega(t)=n^{-1}\sum_i\omega_i(t)$,
% define
% \begin{align}
%  \sigma_\omega(t)&=\sqrt{\frac1n\sum_i[\omega_i(t)-\bar\omega(t)]^2},\nonumber\\
%  s(t)&=10^6\sigma_\omega(t)/\omega_{\rm unit},\qquad\omega_{\rm unit}=1.
%  \label{eq:eval_disagreement_amplitude}
% \end{align}
% Thus $s$ is an instantaneous across-node root-mean-square (RMS)
% amplitude in parts per million (ppm). Its normalization is the fixed
% unit frequency, not $\omega_{\rm base}$ or the reference ramp frequency
% used in Section~\ref{sec:btd_reference}. For $W=[t_s,t_e)$ with duration
% $T_W=t_e-t_s$, report
% \begin{align}
%  R_W&=\sqrt{T_W^{-1}\int_W s(t)^2\,dt},\nonumber\\
%  I_W&=\int_W s(t)\,dt,\qquad P_W=\sup_{t\in W}s(t).
%  \label{eq:eval_frequency_metrics}
% \end{align}
% The recorded integral absolute error (IAE) is $I_W$, the integral of
% this nonnegative RMS amplitude. It is not the integral of the mean
% absolute node deviation. Moments are integrated over complete held-
% frequency intervals using physical-duration weights, not uniform
% averages of controller updates.
% 
% These quantities must also be distinguished from the pairwise spread
% $\Delta_\omega(t)=\max_i\omega_i(t)-\min_i\omega_i(t)$ of
% Section~\ref{sec:btd_reference}. Cauchy--Schwarz gives
% $\Delta_\omega(t)\le\sqrt{2n}\,\sigma_\omega(t)$, but a window RMS
% does not bound a pointwise transient by itself. We retain positive-
% frequency checks and common-frequency diagnostics separately. For M2,
% a supplementary descriptive spread calculation uses the already saved
% shared-time frequency intervals; it is not a new simulation or a
% preregistered primary endpoint.
% 
% \paragraph{Buffer excursion and control activity.}
% Measured buffer excursion and remaining headroom are
% \begin{align}
%  B_W^{\rm obs}&=\max_e\sup_{t\in[t_s,t_e]}|\beta_e(t)-\beta_e^*|,\nonumber\\
%  H_W^{\rm obs}&=\min_e\inf_{t\in[t_s,t_e]}
%                      \min\{\beta_e(t),C-\beta_e(t)\}.
%  \label{eq:eval_buffer_metrics}
% \end{align}
% Every integer phase crossing and ordered arrival is reconstructed.
% Simultaneous events are grouped under the frozen convention, and
% buffer extrema include the right-endpoint state. Display samples are
% not the source of peak or moment statistics. Conservative reference
% certificates are not substituted for measured extrema.
% Mean per-node command total variation (TV) is
% \begin{equation}
%  \mathrm{TV}_W=\frac1n\sum_i
%        \sum_{k:\,t_s\le t^a_{i,k}\le t_e}|u_{i,k}-u_{i,k-1}|.
%  \label{eq:eval_tv}
% \end{equation}
% Activity uses inclusive endpoints in the frozen protocol; adjacent
% subwindow activities must not be added without handling their shared
% boundary. All evaluated M2/M3 runs have zero recorded saturation, so
% requested and applied command increments coincide. TV is a command-
% activity measure, not measured hardware energy or computational cost.
% 
% \paragraph{Windows and recovery.}
% The primary disturbance window is $[2\times10^6,4\times10^6)$.
% M2 additionally separates burst ON and post-OFF intervals at
% $t_{\rm off}=2.1\times10^6$. Before onset, both M2 controllers and all
% M3 cases pass three consecutive $10^5$-tick nominal windows: disagreement
% RMS at most 2 ppm, each buffer's absolute mean bias at most two frames
% and standard deviation at most one frame, common-frequency temporal
% standard deviation at most 1 ppm and three-window mean range at most
% 1 ppm, with positive frequencies/progress and no intervention. These
% are finite engineering centering criteria, not exact equilibrium.
% 
% M2 component recovery requires these criteria in three consecutive
% post-OFF windows. Strict recovery additionally requires the window-mean common
% frequency to return within 2 ppm of that controller's own clean
% trajectory over the same physical-time window. A missing qualifying return by
% $4\times10^6$ is right-censored; it is neither zero recovery time nor
% proof of permanent nonrecovery. M3 uses $[3.6\times10^6,4\times10^6)$
% for tail diagnostics. The separate M2 memory decomposition below uses
% $[3.7\times10^6,4\times10^6)$.
% 
% \subsection{Implementation equivalence and reference consistency}
% \label{subsec:eval_implementation_reference}
% 
% M0 varies polling over 100 and 1000 local ticks, initial frequency
% spread over 35 and 400 ppm, and includes capacity/target settings
% $(2048,1024)$, $(512,128)$ and $(512,384)$ across its six matched pairs.
% All pairs use constant path delay 20, execution delay 10 and a
% $2\times10^5$-tick horizon. Target-centered percentages are used for
% non-midpoint targets. Across 19,800 paired updates, the maximum
% requested/applied correction difference is $5.29\times10^{-18}$.
% Physical and data phases coincide at all evaluated union knots, with
% identical measurement vectors and saved-time buffer counts. No guard
% intervention or saturation is recorded. This supports the implemented
% calibration, integration order and constructor treatment of
% Proposition~\ref{prop:fpi_equivalence}; it is not a nonlinear-Fuzzy
% performance result or a general runtime-equivalence proof.
% 
% Separate previously completed reference calculations call the original
% 25-rule Fuzzy update directly, rather than advance the plant with a
% sector surrogate. Two trajectories use moderate and stronger gain
% coordinates, each with 2000 post-bootstrap updates. For the moderate
% configuration corresponding to Section~\ref{sec:btd_reference}, the
% observed maximum fluid buffer deviation is approximately 17.314056818
% frames, maximum $|X|+|Y|$ is 0.090130790, and maximum absolute correction
% is $3.21398377\times10^{-4}$. There are no observed stepwise
% contraction-envelope violations; frequencies remain in
% $[0.9998,1.0002]$ including startup. These numerical checks are
% consistent with the reference certificate, but the all-time guarantee
% and strict PI comparison come from the interval-prefix and certified-
% tail argument of Theorem~\ref{thm:btd_headroom}, not from a finite
% floating-point trajectory. The reference marking is 1024; the delayed
% actual runs' marking is 1044.
% 
% \subsection{Response to a post-centering latency burst}
% \label{subsec:eval_burst}
% 
% Table~\ref{tab:eval_burst} reports the aligned moderate-gain comparison.
% Figure~\ref{fig:eval_burst_overview} shows the shared-time frequency,
% buffer and latency traces, including the clean counterfactuals.
% Fuzzy reduces RMS-amplitude IAE by 1.864\% and mean command TV by
% 15.938\%. Its disagreement RMS is 0.096\% higher, instantaneous RMS
% peak is 2.396\% higher, and target-referenced buffer excursion is one
% frame larger. These endpoints establish a tradeoff at this deterministic
% operating point, not a uniform disturbance-rejection advantage.
% 
% \begin{table*}[t]
% \caption{M2 metrics over $[2\times10^6,4\times10^6)$ unless stated otherwise.
% Lower is preferable except for buffer headroom. The spread row is a
% supplementary descriptive calculation; no engineering spread threshold
% was predeclared for these runs.}
% \label{tab:eval_burst}
% \centering
% \begin{tabular}{@{}lrr@{}}\toprule
% Metric & Moderate PI & Full Fuzzy\\\midrule
% Disagreement RMS $R_W$ / ppm & 9.67059 & 9.67988\\
% Instantaneous RMS peak $P_W$ / ppm & 32.59447 & 33.37534\\
% RMS-amplitude IAE $I_W/(10^6\ \mathrm{ppm\cdot ticks})$ & 9.94254 & 9.75722\\
% Supplementary peak pairwise spread / ppm & 69.14830 & 70.89161\\
% Buffer excursion $B_W^{\rm obs}$ / frames & 20 & 21\\
% Minimum lower / upper buffer headroom / frames & 1004 / 1017 & 1003 / 1017\\
% Mean command TV / $10^{-4}$ & 2.62204 & 2.20415\\
% Largest command increment / $10^{-5}$ & 4.10000 & 4.21763\\
% Peak common-frequency difference from own clean / ppm & 60.61864 & 61.26423\\
% Component entry / confirmation after OFF / $10^6$ ticks & 0.8 / 1.1 & 0.8 / 1.1\\
% Strict own-clean recovery by the horizon & Censored & Censored\\
% Full-run buffer excursion, including startup / frames & 24 & 21\\\bottomrule
% \end{tabular}
% \end{table*}
% 
% \begin{figure*}[t]
% \centering
% \includegraphics[width=0.97\textwidth]{images/evaluation_20261103/m2_shared_time_response.pdf}
% \caption{Aligned post-centering burst experiment. Continuous-phase Callisto
% uses the fixed moderate PI coordinates and their full-Fuzzy calibration.
% Both controllers have their own clean counterfactual. Desired latency on
% one directed channel changes from 20 to 40 during
% $[2\times10^6,2.1\times10^6)$. Frequency histories use held corrections;
% buffer overview samples are for display, while reported buffer statistics
% come from the complete reconstructed event stream.}
% \Description{Shared-time overview of node frequencies, six receiving-buffer
% occupancies, and the imposed path-latency burst for PI and Fuzzy PI.}
% \label{fig:eval_burst_overview}
% \end{figure*}
% 
% Both controllers remain within capacity with positive frequency and
% continuing logical phase progress, without recorded guard or saturation
% intervention. Including startup, the full-run maximum deviation is
% 24 frames for PI and 21 for Fuzzy. This is a favorable Fuzzy endpoint
% for that full window, whereas the post-onset peaks are 20 and 21.
% Changing the window changes the ranking, so neither figure is used as
% a replacement for the other. Capacities are far from exhaustion.
% 
% Both first enter the window-qualified component recovery band at
% $2.9\times10^6$ and confirm it at $3.2\times10^6$. Thus the experiment
% does not show faster Fuzzy settling under this criterion. Strict
% own-clean recovery is censored for both. In the memory-diagnostic
% tail, the common-frequency changes relative to own clean are
% $-59.40198$ ppm for PI and $-60.51069$ ppm for Fuzzy. A stationary
% common frequency, frequency agreement between nodes, return to own
% clean, and return to frequency one are different properties.
% 
% \subsection{Observed nonlinear mechanism and the first buffer peak}
% \label{subsec:eval_mechanism}
% 
% The M2 rule replay agrees with the full nonlinear implementation.
% All sampled Fuzzy input pairs remain in $\Omega_5$, with disturbance-window
% maximum $|X|+|Y|=2.001888$. The central identity of
% Proposition~\ref{prop:fpi_central} therefore applies to those observed
% updates. However, this trajectory does not stay in the smaller
% $\Omega_1$ input diamond used by the reference example. This is not
% evidence of a sector violation or instability; it is another reason,
% besides asynchrony, delay and quantization, that the example's
% invariant-region certificate is not a certificate for this actual burst.
% 
% The Section~\ref{sec:fpi_controller} decomposition into proportional,
% integral-like and accumulated nonlinear-residual terms reproduces the
% requested correction with maximum residual approximately
% $6.46\times10^{-18}$. In the diagnostic tail, the Fuzzy own-clean
% common-frequency shift decomposes into
% \begin{equation}
%  -60.51069=-57.61859-2.89210\quad\mathrm{ppm},
%  \label{eq:eval_memory_accounting}
% \end{equation}
% where the two terms are the linear integral-like contribution and
% nonlinear memory; the proportional contribution is zero. PI's
% $-59.40198$ ppm shift is its integral contribution. This is exact
% same-trajectory accounting up to numerical arithmetic, not an
% independently evolving $F_0$ closed-loop ablation or a unique causal
% explanation of the physical common frequency.
% 
% Figure~\ref{fig:eval_mechanism} also resolves the first buffer peak.
% PI and Fuzzy first attain their disturbance-window maximum target deficits at
% ON$+39.105$ and ON$+40.731$ ticks, whereas their first changed commands
% relative to own clean occur at ON$+155.216$ and ON$+155.841$.
% At the respective first extrema, clean/burst occupancies are
% 1024/1004 for PI and 1023/1003 for Fuzzy. Both same-time
% clean-minus-burst deficits are therefore 20 frames: target deficits of
% 20 and 21 combine that delivery deficit with clean-state offsets of
% zero and one. Consequently the first attainment is not a measurement
% of post-burst feedback effectiveness. The minimum is revisited later;
% we do not claim that every peak instance precedes feedback.
% 
% \begin{figure*}[t]
% \centering
% \includegraphics[width=0.97\textwidth]{images/evaluation_20261103/m2_nonlinear_memory.pdf}
% \caption{Nonlinear mechanism diagnostics reusing the four M2 runs.
% The panels show the exact central residual $F-F_0$, common-frequency
% memory accounting, command activity and complete early buffer events.
% The hatched $F_0$ activity bar is a replay on Fuzzy's measured inputs,
% not a third closed-loop experiment. The residual surface is multiplied
% by negative $K_u$ in the controller. The memory panel uses
% $[3.7\times10^6,4\times10^6)$; activity uses the main disturbance
% window. First peak attainment precedes the first changed command, but
% the peak can be revisited after feedback begins.}
% \Description{Four panels connecting Fuzzy rule residuals, stored correction
% memory, command total variation, and early buffer depletion timing.}
% \label{fig:eval_mechanism}
% \end{figure*}
% 
% \subsection{Independent paired random-latency validation}
% \label{subsec:eval_random}
% 
% All 180 M3 validation runs are normal-admissible, with zero saved-state
% mismatches in complete-event replay. Table~\ref{tab:eval_random}
% summarizes the primary endpoint. The ratio estimand is
% \begin{equation}
%  \widehat r=\frac1{30}\sum_{s=1}^{30}
%        \frac{R_{W,s}^{\rm Fuzzy}}{R_{W,s}^{\rm PI}},
%  \label{eq:eval_paired_ratio}
% \end{equation}
% where $s$ indexes independent seeds. It is not the ratio of the two
% group means. Pointwise 95\% confidence intervals (CIs) use 20,000
% paired-seed percentile-bootstrap resamples with fixed resampling seeds.
% Time points, edges and bootstrap draws are not independent experimental
% replications. Regimes are analyzed separately; the intervals do not
% provide simultaneous multiple-endpoint coverage. There are no excluded
% pairs in the reported comparisons.
% 
% \begin{table*}[t]
% \caption{M3 primary disturbance-window disagreement RMS. Controller entries
% are means over 30 seeds in ppm. Ratios are means of paired seed ratios
% with pointwise 95\% CIs. A ratio below one favors Fuzzy.}
% \label{tab:eval_random}
% \centering
% \begin{tabular}{@{}lrrrll@{}}\toprule
% Profile grid & PI moderate & PI stronger & Fuzzy & Fuzzy/moderate [CI] & Fuzzy/stronger [CI]\\\midrule
% $T_c=1000$ & 4.911 & 9.347 & 5.856 & 1.1925 [1.1722,1.2150] & 0.6266 [0.6155,0.6389]\\
% $T_c=10000$ & 6.696 & 9.699 & 7.152 & 1.0679 [1.0624,1.0736] & 0.7365 [0.7184,0.7554]\\\bottomrule
% \end{tabular}
% \end{table*}
% 
% \begin{figure*}[t]
% \centering
% \includegraphics[width=0.97\textwidth]{images/evaluation_20261103/m3_primary_paired_effects.pdf}
% \caption{Independent paired random-latency validation: primary disagreement
% RMS and paired Fuzzy/PI effects for two physical profile grids. Thirty
% seeds per regime are shared across the three fixed controllers. The
% moderate comparator is the primary axis-matched PI baseline; stronger
% PI is a separate fixed candidate. Ratios and pointwise confidence
% intervals follow the paired-seed protocol, rather than resampling
% correlated controller updates.}
% \Description{Primary random-latency disagreement results and paired effects
% against moderate and stronger fixed PI for thirty seeds per regime.}
% \label{fig:eval_random_primary}
% \end{figure*}
% 
% Relative to moderate PI, Fuzzy has higher RMS in all 30 seeds of each
% regime: mean paired increases are 19.245\% and 6.791\%. The lower
% RMS relative to stronger fixed PI is comparator-specific; moderate PI
% already has lower RMS than Fuzzy, so that comparison does not establish
% an inherent advantage of fuzzy inference over suitable PI tuning.
% The frozen random-latency hypothesis of lower disagreement RMS is
% therefore not supported against the primary baseline in these regimes.
% 
% Within the disturbance windows, all runs have at most seven frames of
% actual target deviation and at least 1017 frames of lower/upper
% headroom. These extrema exclude startup. They establish finite-run
% admissibility in this interior operating region, not a larger Fuzzy
% safe set. The random result also does not extend the M2 TV reduction:
% mean paired Fuzzy/moderate TV ratios are 1.03045 and 1.03153.
% Figure~\ref{fig:eval_random_secondary} retains buffer, effort and
% common-frequency diagnostics. Mean tail common-frequency shifts from
% each controller's pre-disturbance level are approximately
% $-156.682$ versus $-72.576$ ppm at $T_c=1000$, and $-76.540$
% versus $-67.816$ ppm at $T_c=10000$ (Fuzzy versus moderate PI).
% These finite-horizon shifts and trends are reported separately from
% relative synchronization; no absolute-frequency loop is present.
% 
% \begin{figure*}[t]
% \centering
% \includegraphics[width=0.97\textwidth]{images/evaluation_20261103/m3_buffer_effort_drift.pdf}
% \caption{Secondary M3 diagnostics over the declared windows: complete-event
% buffer excursions, command activity and common-frequency shift.
% The same 30 paired seeds underlie each regime; additional endpoints
% do not create additional independent replications. Buffer extrema
% refer to the disturbance window, not the startup trajectory.}
% \Description{Buffer, control activity and common-frequency diagnostics
% complementing the primary random-latency synchronization metric.}
% \label{fig:eval_random_secondary}
% \end{figure*}
% 
% \subsection{What the completed evaluation establishes}
% \label{subsec:eval_support}
% 
% The evidence supports the implementation-consistent PI calibration,
% the exact nonlinear mechanism on observed M2 inputs, and the
% conditional reference-model buffer-headroom theorem. Actual Callisto
% demonstrates endpoint-dependent response shaping: a smaller full-run
% buffer excursion and lower post-burst IAE/TV at the stated M2 point,
% alongside higher post-onset peaks and no faster component recovery.
% The independent M3 study limits any general random-jitter advantage
% claim for the frozen controller. These outcomes are consistent with
% the paper's conditional mechanism-and-tradeoff formulation; they do
% not support uniform Fuzzy superiority.
% 
% The completed actual campaigns do not test the boundary of the
% reference invariant set. Their large buffer reserves cannot substitute
% for certified-safe, uncertified-but-observed-safe and failed cases
% around a common boundary. Nor do finite observations prove
% infinite-horizon safety or the reference theorem's optional
% frequency-spread tolerance in asynchronous delayed execution.
% Observed phase progress and capacity compliance support normal model
% execution; application-payload correctness and a formal refinement
% proof of logical synchrony are outside this evaluation.
% 
% Finally, these are fixed-controller comparisons, not a completed
% equal-budget tuning/held-out-test study. A separate preregistered
% tuning campaign has not completed its audit, selection and untouched
% tests; its partial results are excluded from the present performance
% tables. The evaluation also provides no P/GABA comparison, switched
% FFP or observer result, dropout recovery, serialized-link hardware
% test, or meaningful controller runtime/energy ranking. Consequently,
% practical adoption beyond the reported operating points remains an
% empirical question even though the reference-model headroom benefit
% is mathematically established.

\section{Discussion and Limitations}
\label{sec:discussion}
\noindent\textbf{Draft status: to be written.}
The full discussion remains pending. The certificate in
Section~\ref{sec:btd_reference} applies to its specified synchronous,
zero-delay fluid reference and prescribed initialization. It does not
certify asynchronous, delayed, quantized Callisto execution or establish
universal superiority over PI.

\section{Conclusion}
\label{sec:conclusion}
\noindent\textbf{Draft status: to be written.}
Complete this section after the evaluation and discussion are integrated,
keeping reference-model guarantees distinct from measured Callisto results.


\begin{acks}
\textbf{Draft status.} Funding and acknowledgments are pending. Remove
identifying material when preparing an anonymous submission.
\end{acks}

\bibliographystyle{ACM-Reference-Format}
\bibliography{references}

\end{document}

% Historical draft abstract placeholder.
% 
% \textbf{Draft status.} The abstract is pending. This working manuscript
% integrates the Introduction, Background, expanded controller model,
% reference-model derivations and Evaluation; Discussion and Conclusion
% remain to be written. Author, conference and CCS metadata require final confirmation.
